\subsection{Security, Safety, and Governance}
\paragraph{MCQ-1721 [medium]} Which statement best describes Prompt injection?
\begin{enumerate}[label=\Alph*.]
\item Adversarial testing explores misuse, prompt injection, privacy, bias, jailbreaks, tool abuse, and failure recovery.
\item Weights, datasets, containers, packages, adapters, and prompts require provenance, scanning, signing, and review.
\item Untrusted content attempts to override instructions, exfiltrate data, or induce tool actions.
\item Immutable event records connect identity, input, policy, model, evidence, tool calls, decisions, and outcomes.
\end{enumerate}
\noindent\textbf{Answer.} Untrusted content attempts to override instructions, exfiltrate data, or induce tool actions.
\par\textit{Untrusted content attempts to override instructions, exfiltrate data, or induce tool actions. Tradeoff: Content isolation and least privilege reduce impact but cannot make untrusted text authoritative. Watch for: Telling the model to ignore attacks is not a security boundary.}
\paragraph{MCQ-1722 [hard]} What is the central engineering tradeoff of Prompt injection?
\begin{enumerate}[label=\Alph*.]
\item Content isolation and least privilege reduce impact but cannot make untrusted text authoritative.
\item Investigations and compliance improve, while logs themselves become sensitive assets.
\item It reveals unknown weaknesses but provides samples, not a guarantee of safety.
\item Assurance improves while promotion pipelines become more deliberate.
\end{enumerate}
\noindent\textbf{Answer.} Content isolation and least privilege reduce impact but cannot make untrusted text authoritative.
\par\textit{Untrusted content attempts to override instructions, exfiltrate data, or induce tool actions. Tradeoff: Content isolation and least privilege reduce impact but cannot make untrusted text authoritative. Watch for: Telling the model to ignore attacks is not a security boundary.}
\paragraph{MCQ-1723 [hard]} Which failure mode should an interviewer expect you to identify for Prompt injection?
\begin{enumerate}[label=\Alph*.]
\item Running one generic jailbreak list misses application-specific assets and actions.
\item Telling the model to ignore attacks is not a security boundary.
\item Loading arbitrary pickle-based artifacts can execute code.
\item Logging reasoning text instead of decisive actions and evidence can reduce clarity.
\end{enumerate}
\noindent\textbf{Answer.} Telling the model to ignore attacks is not a security boundary.
\par\textit{Untrusted content attempts to override instructions, exfiltrate data, or induce tool actions. Tradeoff: Content isolation and least privilege reduce impact but cannot make untrusted text authoritative. Watch for: Telling the model to ignore attacks is not a security boundary.}
\paragraph{MCQ-1724 [medium]} A design review is specifically concerned with this risk: Telling the model to ignore attacks is not a security boundary. Which topic should be reviewed first?
\begin{enumerate}[label=\Alph*.]
\item Auditability
\item Prompt injection
\item Red teaming
\item Model supply chain
\end{enumerate}
\noindent\textbf{Answer.} Prompt injection
\par\textit{Untrusted content attempts to override instructions, exfiltrate data, or induce tool actions. Tradeoff: Content isolation and least privilege reduce impact but cannot make untrusted text authoritative. Watch for: Telling the model to ignore attacks is not a security boundary.}
\paragraph{MCQ-1725 [hard]} Which design note most accurately balances the benefits and costs of Prompt injection?
\begin{enumerate}[label=\Alph*.]
\item Content isolation and least privilege reduce impact but cannot make untrusted text authoritative.
\item It reveals unknown weaknesses but provides samples, not a guarantee of safety.
\item Investigations and compliance improve, while logs themselves become sensitive assets.
\item Assurance improves while promotion pipelines become more deliberate.
\end{enumerate}
\noindent\textbf{Answer.} Content isolation and least privilege reduce impact but cannot make untrusted text authoritative.
\par\textit{Untrusted content attempts to override instructions, exfiltrate data, or induce tool actions. Tradeoff: Content isolation and least privilege reduce impact but cannot make untrusted text authoritative. Watch for: Telling the model to ignore attacks is not a security boundary.}
\paragraph{MCQ-1726 [medium]} Which explanation would be strongest in a system-design interview about Prompt injection?
\begin{enumerate}[label=\Alph*.]
\item Untrusted content attempts to override instructions, exfiltrate data, or induce tool actions.
\item Adversarial testing explores misuse, prompt injection, privacy, bias, jailbreaks, tool abuse, and failure recovery.
\item Weights, datasets, containers, packages, adapters, and prompts require provenance, scanning, signing, and review.
\item Immutable event records connect identity, input, policy, model, evidence, tool calls, decisions, and outcomes.
\end{enumerate}
\noindent\textbf{Answer.} Untrusted content attempts to override instructions, exfiltrate data, or induce tool actions.
\par\textit{Untrusted content attempts to override instructions, exfiltrate data, or induce tool actions. Tradeoff: Content isolation and least privilege reduce impact but cannot make untrusted text authoritative. Watch for: Telling the model to ignore attacks is not a security boundary.}
\paragraph{MCQ-1727 [hard]} Which production warning is most directly associated with Prompt injection?
\begin{enumerate}[label=\Alph*.]
\item Running one generic jailbreak list misses application-specific assets and actions.
\item Loading arbitrary pickle-based artifacts can execute code.
\item Telling the model to ignore attacks is not a security boundary.
\item Logging reasoning text instead of decisive actions and evidence can reduce clarity.
\end{enumerate}
\noindent\textbf{Answer.} Telling the model to ignore attacks is not a security boundary.
\par\textit{Untrusted content attempts to override instructions, exfiltrate data, or induce tool actions. Tradeoff: Content isolation and least privilege reduce impact but cannot make untrusted text authoritative. Watch for: Telling the model to ignore attacks is not a security boundary.}
\paragraph{FC-1728 [medium]} Define Prompt injection in interview-ready language.
\noindent\textbf{Answer.} Untrusted content attempts to override instructions, exfiltrate data, or induce tool actions.
\par\textit{Untrusted content attempts to override instructions, exfiltrate data, or induce tool actions.}
\paragraph{FC-1729 [hard]} State the main tradeoff and production pitfall for Prompt injection.
\noindent\textbf{Answer.} Tradeoff: Content isolation and least privilege reduce impact but cannot make untrusted text authoritative. Pitfall: Telling the model to ignore attacks is not a security boundary.
\par\textit{Tradeoff: Content isolation and least privilege reduce impact but cannot make untrusted text authoritative. Pitfall: Telling the model to ignore attacks is not a security boundary.}
\paragraph{LA-1730 [hard]} Explain Prompt injection from first principles, then show how you would evaluate and operate it in production.
\noindent\textbf{Answer.} Start with the mechanism: Untrusted content attempts to override instructions, exfiltrate data, or induce tool actions. Make the tradeoff explicit: Content isolation and least privilege reduce impact but cannot make untrusted text authoritative. Define workload-specific offline metrics and a latency/cost budget, test important slices, instrument the critical path, and create a rollback criterion. The production review must explicitly guard against this failure: Telling the model to ignore attacks is not a security boundary.
\par\textit{A strong answer connects mechanism, measurable tradeoffs, failure isolation, observability, and rollback rather than listing product features.}
\paragraph{MCQ-1731 [medium]} Which statement best describes Tool authorization?
\begin{enumerate}[label=\Alph*.]
\item Untrusted content attempts to override instructions, exfiltrate data, or induce tool actions.
\item Credentials are issued narrowly, stored outside prompts and code, rotated, audited, and redacted from telemetry.
\item Risk processes map context, measure impacts, manage controls, document evidence, and govern residual risk.
\item The host checks identity, scope, resource, action, policy, and user intent before every tool execution.
\end{enumerate}
\noindent\textbf{Answer.} The host checks identity, scope, resource, action, policy, and user intent before every tool execution.
\par\textit{The host checks identity, scope, resource, action, policy, and user intent before every tool execution. Tradeoff: Fine-grained authorization limits blast radius but adds policy complexity. Watch for: Authorizing the agent once at session start enables confused-deputy attacks.}
\paragraph{MCQ-1732 [hard]} What is the central engineering tradeoff of Tool authorization?
\begin{enumerate}[label=\Alph*.]
\item Governance improves accountability but can become checkbox compliance without technical signals.
\item Fine-grained authorization limits blast radius but adds policy complexity.
\item Short-lived credentials reduce exposure but require reliable identity infrastructure.
\item Content isolation and least privilege reduce impact but cannot make untrusted text authoritative.
\end{enumerate}
\noindent\textbf{Answer.} Fine-grained authorization limits blast radius but adds policy complexity.
\par\textit{The host checks identity, scope, resource, action, policy, and user intent before every tool execution. Tradeoff: Fine-grained authorization limits blast radius but adds policy complexity. Watch for: Authorizing the agent once at session start enables confused-deputy attacks.}
\paragraph{MCQ-1733 [hard]} Which failure mode should an interviewer expect you to identify for Tool authorization?
\begin{enumerate}[label=\Alph*.]
\item Authorizing the agent once at session start enables confused-deputy attacks.
\item Telling the model to ignore attacks is not a security boundary.
\item Environment variables copied into sandboxes or crash logs leak broad authority.
\item A model card without deployment-specific threat modeling is incomplete.
\end{enumerate}
\noindent\textbf{Answer.} Authorizing the agent once at session start enables confused-deputy attacks.
\par\textit{The host checks identity, scope, resource, action, policy, and user intent before every tool execution. Tradeoff: Fine-grained authorization limits blast radius but adds policy complexity. Watch for: Authorizing the agent once at session start enables confused-deputy attacks.}
\paragraph{MCQ-1734 [medium]} A design review is specifically concerned with this risk: Authorizing the agent once at session start enables confused-deputy attacks. Which topic should be reviewed first?
\begin{enumerate}[label=\Alph*.]
\item AI risk management
\item Secrets management
\item Tool authorization
\item Prompt injection
\end{enumerate}
\noindent\textbf{Answer.} Tool authorization
\par\textit{The host checks identity, scope, resource, action, policy, and user intent before every tool execution. Tradeoff: Fine-grained authorization limits blast radius but adds policy complexity. Watch for: Authorizing the agent once at session start enables confused-deputy attacks.}
\paragraph{MCQ-1735 [hard]} Which design note most accurately balances the benefits and costs of Tool authorization?
\begin{enumerate}[label=\Alph*.]
\item Content isolation and least privilege reduce impact but cannot make untrusted text authoritative.
\item Fine-grained authorization limits blast radius but adds policy complexity.
\item Governance improves accountability but can become checkbox compliance without technical signals.
\item Short-lived credentials reduce exposure but require reliable identity infrastructure.
\end{enumerate}
\noindent\textbf{Answer.} Fine-grained authorization limits blast radius but adds policy complexity.
\par\textit{The host checks identity, scope, resource, action, policy, and user intent before every tool execution. Tradeoff: Fine-grained authorization limits blast radius but adds policy complexity. Watch for: Authorizing the agent once at session start enables confused-deputy attacks.}
\paragraph{MCQ-1736 [medium]} Which explanation would be strongest in a system-design interview about Tool authorization?
\begin{enumerate}[label=\Alph*.]
\item Untrusted content attempts to override instructions, exfiltrate data, or induce tool actions.
\item The host checks identity, scope, resource, action, policy, and user intent before every tool execution.
\item Credentials are issued narrowly, stored outside prompts and code, rotated, audited, and redacted from telemetry.
\item Risk processes map context, measure impacts, manage controls, document evidence, and govern residual risk.
\end{enumerate}
\noindent\textbf{Answer.} The host checks identity, scope, resource, action, policy, and user intent before every tool execution.
\par\textit{The host checks identity, scope, resource, action, policy, and user intent before every tool execution. Tradeoff: Fine-grained authorization limits blast radius but adds policy complexity. Watch for: Authorizing the agent once at session start enables confused-deputy attacks.}
\paragraph{MCQ-1737 [hard]} Which production warning is most directly associated with Tool authorization?
\begin{enumerate}[label=\Alph*.]
\item Authorizing the agent once at session start enables confused-deputy attacks.
\item Telling the model to ignore attacks is not a security boundary.
\item A model card without deployment-specific threat modeling is incomplete.
\item Environment variables copied into sandboxes or crash logs leak broad authority.
\end{enumerate}
\noindent\textbf{Answer.} Authorizing the agent once at session start enables confused-deputy attacks.
\par\textit{The host checks identity, scope, resource, action, policy, and user intent before every tool execution. Tradeoff: Fine-grained authorization limits blast radius but adds policy complexity. Watch for: Authorizing the agent once at session start enables confused-deputy attacks.}
\paragraph{FC-1738 [medium]} Define Tool authorization in interview-ready language.
\noindent\textbf{Answer.} The host checks identity, scope, resource, action, policy, and user intent before every tool execution.
\par\textit{The host checks identity, scope, resource, action, policy, and user intent before every tool execution.}
\paragraph{FC-1739 [hard]} State the main tradeoff and production pitfall for Tool authorization.
\noindent\textbf{Answer.} Tradeoff: Fine-grained authorization limits blast radius but adds policy complexity. Pitfall: Authorizing the agent once at session start enables confused-deputy attacks.
\par\textit{Tradeoff: Fine-grained authorization limits blast radius but adds policy complexity. Pitfall: Authorizing the agent once at session start enables confused-deputy attacks.}
\paragraph{LA-1740 [hard]} Explain Tool authorization from first principles, then show how you would evaluate and operate it in production.
\noindent\textbf{Answer.} Start with the mechanism: The host checks identity, scope, resource, action, policy, and user intent before every tool execution. Make the tradeoff explicit: Fine-grained authorization limits blast radius but adds policy complexity. Define workload-specific offline metrics and a latency/cost budget, test important slices, instrument the critical path, and create a rollback criterion. The production review must explicitly guard against this failure: Authorizing the agent once at session start enables confused-deputy attacks.
\par\textit{A strong answer connects mechanism, measurable tradeoffs, failure isolation, observability, and rollback rather than listing product features.}
\paragraph{MCQ-1741 [medium]} Which statement best describes Secrets management?
\begin{enumerate}[label=\Alph*.]
\item Weights, datasets, containers, packages, adapters, and prompts require provenance, scanning, signing, and review.
\item Structured schema, semantic checks, policy filters, allowlists, and deterministic verification gate model output.
\item Data minimization, purpose limitation, access control, retention, deletion, and regional handling apply to prompts and traces.
\item Credentials are issued narrowly, stored outside prompts and code, rotated, audited, and redacted from telemetry.
\end{enumerate}
\noindent\textbf{Answer.} Credentials are issued narrowly, stored outside prompts and code, rotated, audited, and redacted from telemetry.
\par\textit{Credentials are issued narrowly, stored outside prompts and code, rotated, audited, and redacted from telemetry. Tradeoff: Short-lived credentials reduce exposure but require reliable identity infrastructure. Watch for: Environment variables copied into sandboxes or crash logs leak broad authority.}
\paragraph{MCQ-1742 [hard]} What is the central engineering tradeoff of Secrets management?
\begin{enumerate}[label=\Alph*.]
\item Short-lived credentials reduce exposure but require reliable identity infrastructure.
\item Privacy protection can reduce debugging data and personalization.
\item Fail-closed validation improves safety but may lower completion rate.
\item Assurance improves while promotion pipelines become more deliberate.
\end{enumerate}
\noindent\textbf{Answer.} Short-lived credentials reduce exposure but require reliable identity infrastructure.
\par\textit{Credentials are issued narrowly, stored outside prompts and code, rotated, audited, and redacted from telemetry. Tradeoff: Short-lived credentials reduce exposure but require reliable identity infrastructure. Watch for: Environment variables copied into sandboxes or crash logs leak broad authority.}
\paragraph{MCQ-1743 [hard]} Which failure mode should an interviewer expect you to identify for Secrets management?
\begin{enumerate}[label=\Alph*.]
\item A valid JSON schema does not prove safe SQL, URLs, or business intent.
\item Redacting only user messages while retaining retrieved documents leaves exposure.
\item Loading arbitrary pickle-based artifacts can execute code.
\item Environment variables copied into sandboxes or crash logs leak broad authority.
\end{enumerate}
\noindent\textbf{Answer.} Environment variables copied into sandboxes or crash logs leak broad authority.
\par\textit{Credentials are issued narrowly, stored outside prompts and code, rotated, audited, and redacted from telemetry. Tradeoff: Short-lived credentials reduce exposure but require reliable identity infrastructure. Watch for: Environment variables copied into sandboxes or crash logs leak broad authority.}
\paragraph{MCQ-1744 [medium]} A design review is specifically concerned with this risk: Environment variables copied into sandboxes or crash logs leak broad authority. Which topic should be reviewed first?
\begin{enumerate}[label=\Alph*.]
\item Secrets management
\item Model supply chain
\item Output validation
\item Data privacy
\end{enumerate}
\noindent\textbf{Answer.} Secrets management
\par\textit{Credentials are issued narrowly, stored outside prompts and code, rotated, audited, and redacted from telemetry. Tradeoff: Short-lived credentials reduce exposure but require reliable identity infrastructure. Watch for: Environment variables copied into sandboxes or crash logs leak broad authority.}
\paragraph{MCQ-1745 [hard]} Which design note most accurately balances the benefits and costs of Secrets management?
\begin{enumerate}[label=\Alph*.]
\item Short-lived credentials reduce exposure but require reliable identity infrastructure.
\item Privacy protection can reduce debugging data and personalization.
\item Fail-closed validation improves safety but may lower completion rate.
\item Assurance improves while promotion pipelines become more deliberate.
\end{enumerate}
\noindent\textbf{Answer.} Short-lived credentials reduce exposure but require reliable identity infrastructure.
\par\textit{Credentials are issued narrowly, stored outside prompts and code, rotated, audited, and redacted from telemetry. Tradeoff: Short-lived credentials reduce exposure but require reliable identity infrastructure. Watch for: Environment variables copied into sandboxes or crash logs leak broad authority.}
\paragraph{MCQ-1746 [medium]} Which explanation would be strongest in a system-design interview about Secrets management?
\begin{enumerate}[label=\Alph*.]
\item Data minimization, purpose limitation, access control, retention, deletion, and regional handling apply to prompts and traces.
\item Credentials are issued narrowly, stored outside prompts and code, rotated, audited, and redacted from telemetry.
\item Structured schema, semantic checks, policy filters, allowlists, and deterministic verification gate model output.
\item Weights, datasets, containers, packages, adapters, and prompts require provenance, scanning, signing, and review.
\end{enumerate}
\noindent\textbf{Answer.} Credentials are issued narrowly, stored outside prompts and code, rotated, audited, and redacted from telemetry.
\par\textit{Credentials are issued narrowly, stored outside prompts and code, rotated, audited, and redacted from telemetry. Tradeoff: Short-lived credentials reduce exposure but require reliable identity infrastructure. Watch for: Environment variables copied into sandboxes or crash logs leak broad authority.}
\paragraph{MCQ-1747 [hard]} Which production warning is most directly associated with Secrets management?
\begin{enumerate}[label=\Alph*.]
\item Redacting only user messages while retaining retrieved documents leaves exposure.
\item Environment variables copied into sandboxes or crash logs leak broad authority.
\item Loading arbitrary pickle-based artifacts can execute code.
\item A valid JSON schema does not prove safe SQL, URLs, or business intent.
\end{enumerate}
\noindent\textbf{Answer.} Environment variables copied into sandboxes or crash logs leak broad authority.
\par\textit{Credentials are issued narrowly, stored outside prompts and code, rotated, audited, and redacted from telemetry. Tradeoff: Short-lived credentials reduce exposure but require reliable identity infrastructure. Watch for: Environment variables copied into sandboxes or crash logs leak broad authority.}
\paragraph{FC-1748 [medium]} Define Secrets management in interview-ready language.
\noindent\textbf{Answer.} Credentials are issued narrowly, stored outside prompts and code, rotated, audited, and redacted from telemetry.
\par\textit{Credentials are issued narrowly, stored outside prompts and code, rotated, audited, and redacted from telemetry.}
\paragraph{FC-1749 [hard]} State the main tradeoff and production pitfall for Secrets management.
\noindent\textbf{Answer.} Tradeoff: Short-lived credentials reduce exposure but require reliable identity infrastructure. Pitfall: Environment variables copied into sandboxes or crash logs leak broad authority.
\par\textit{Tradeoff: Short-lived credentials reduce exposure but require reliable identity infrastructure. Pitfall: Environment variables copied into sandboxes or crash logs leak broad authority.}
\paragraph{LA-1750 [hard]} Explain Secrets management from first principles, then show how you would evaluate and operate it in production.
\noindent\textbf{Answer.} Start with the mechanism: Credentials are issued narrowly, stored outside prompts and code, rotated, audited, and redacted from telemetry. Make the tradeoff explicit: Short-lived credentials reduce exposure but require reliable identity infrastructure. Define workload-specific offline metrics and a latency/cost budget, test important slices, instrument the critical path, and create a rollback criterion. The production review must explicitly guard against this failure: Environment variables copied into sandboxes or crash logs leak broad authority.
\par\textit{A strong answer connects mechanism, measurable tradeoffs, failure isolation, observability, and rollback rather than listing product features.}
\paragraph{MCQ-1751 [medium]} Which statement best describes Sandboxing?
\begin{enumerate}[label=\Alph*.]
\item Credentials are issued narrowly, stored outside prompts and code, rotated, audited, and redacted from telemetry.
\item Immutable event records connect identity, input, policy, model, evidence, tool calls, decisions, and outcomes.
\item The host checks identity, scope, resource, action, policy, and user intent before every tool execution.
\item Defense in depth combines process, filesystem, network, syscall, credential, resource, and time boundaries.
\end{enumerate}
\noindent\textbf{Answer.} Defense in depth combines process, filesystem, network, syscall, credential, resource, and time boundaries.
\par\textit{Defense in depth combines process, filesystem, network, syscall, credential, resource, and time boundaries. Tradeoff: Strong isolation raises operational cost and may restrict legitimate tools. Watch for: Running as non-root inside an otherwise privileged container is insufficient.}
\paragraph{MCQ-1752 [hard]} What is the central engineering tradeoff of Sandboxing?
\begin{enumerate}[label=\Alph*.]
\item Strong isolation raises operational cost and may restrict legitimate tools.
\item Short-lived credentials reduce exposure but require reliable identity infrastructure.
\item Fine-grained authorization limits blast radius but adds policy complexity.
\item Investigations and compliance improve, while logs themselves become sensitive assets.
\end{enumerate}
\noindent\textbf{Answer.} Strong isolation raises operational cost and may restrict legitimate tools.
\par\textit{Defense in depth combines process, filesystem, network, syscall, credential, resource, and time boundaries. Tradeoff: Strong isolation raises operational cost and may restrict legitimate tools. Watch for: Running as non-root inside an otherwise privileged container is insufficient.}
\paragraph{MCQ-1753 [hard]} Which failure mode should an interviewer expect you to identify for Sandboxing?
\begin{enumerate}[label=\Alph*.]
\item Environment variables copied into sandboxes or crash logs leak broad authority.
\item Logging reasoning text instead of decisive actions and evidence can reduce clarity.
\item Running as non-root inside an otherwise privileged container is insufficient.
\item Authorizing the agent once at session start enables confused-deputy attacks.
\end{enumerate}
\noindent\textbf{Answer.} Running as non-root inside an otherwise privileged container is insufficient.
\par\textit{Defense in depth combines process, filesystem, network, syscall, credential, resource, and time boundaries. Tradeoff: Strong isolation raises operational cost and may restrict legitimate tools. Watch for: Running as non-root inside an otherwise privileged container is insufficient.}
\paragraph{MCQ-1754 [medium]} A design review is specifically concerned with this risk: Running as non-root inside an otherwise privileged container is insufficient. Which topic should be reviewed first?
\begin{enumerate}[label=\Alph*.]
\item Auditability
\item Secrets management
\item Tool authorization
\item Sandboxing
\end{enumerate}
\noindent\textbf{Answer.} Sandboxing
\par\textit{Defense in depth combines process, filesystem, network, syscall, credential, resource, and time boundaries. Tradeoff: Strong isolation raises operational cost and may restrict legitimate tools. Watch for: Running as non-root inside an otherwise privileged container is insufficient.}
\paragraph{MCQ-1755 [hard]} Which design note most accurately balances the benefits and costs of Sandboxing?
\begin{enumerate}[label=\Alph*.]
\item Fine-grained authorization limits blast radius but adds policy complexity.
\item Short-lived credentials reduce exposure but require reliable identity infrastructure.
\item Strong isolation raises operational cost and may restrict legitimate tools.
\item Investigations and compliance improve, while logs themselves become sensitive assets.
\end{enumerate}
\noindent\textbf{Answer.} Strong isolation raises operational cost and may restrict legitimate tools.
\par\textit{Defense in depth combines process, filesystem, network, syscall, credential, resource, and time boundaries. Tradeoff: Strong isolation raises operational cost and may restrict legitimate tools. Watch for: Running as non-root inside an otherwise privileged container is insufficient.}
\paragraph{MCQ-1756 [medium]} Which explanation would be strongest in a system-design interview about Sandboxing?
\begin{enumerate}[label=\Alph*.]
\item Immutable event records connect identity, input, policy, model, evidence, tool calls, decisions, and outcomes.
\item Defense in depth combines process, filesystem, network, syscall, credential, resource, and time boundaries.
\item The host checks identity, scope, resource, action, policy, and user intent before every tool execution.
\item Credentials are issued narrowly, stored outside prompts and code, rotated, audited, and redacted from telemetry.
\end{enumerate}
\noindent\textbf{Answer.} Defense in depth combines process, filesystem, network, syscall, credential, resource, and time boundaries.
\par\textit{Defense in depth combines process, filesystem, network, syscall, credential, resource, and time boundaries. Tradeoff: Strong isolation raises operational cost and may restrict legitimate tools. Watch for: Running as non-root inside an otherwise privileged container is insufficient.}
\paragraph{MCQ-1757 [hard]} Which production warning is most directly associated with Sandboxing?
\begin{enumerate}[label=\Alph*.]
\item Running as non-root inside an otherwise privileged container is insufficient.
\item Authorizing the agent once at session start enables confused-deputy attacks.
\item Logging reasoning text instead of decisive actions and evidence can reduce clarity.
\item Environment variables copied into sandboxes or crash logs leak broad authority.
\end{enumerate}
\noindent\textbf{Answer.} Running as non-root inside an otherwise privileged container is insufficient.
\par\textit{Defense in depth combines process, filesystem, network, syscall, credential, resource, and time boundaries. Tradeoff: Strong isolation raises operational cost and may restrict legitimate tools. Watch for: Running as non-root inside an otherwise privileged container is insufficient.}
\paragraph{FC-1758 [medium]} Define Sandboxing in interview-ready language.
\noindent\textbf{Answer.} Defense in depth combines process, filesystem, network, syscall, credential, resource, and time boundaries.
\par\textit{Defense in depth combines process, filesystem, network, syscall, credential, resource, and time boundaries.}
\paragraph{FC-1759 [hard]} State the main tradeoff and production pitfall for Sandboxing.
\noindent\textbf{Answer.} Tradeoff: Strong isolation raises operational cost and may restrict legitimate tools. Pitfall: Running as non-root inside an otherwise privileged container is insufficient.
\par\textit{Tradeoff: Strong isolation raises operational cost and may restrict legitimate tools. Pitfall: Running as non-root inside an otherwise privileged container is insufficient.}
\paragraph{LA-1760 [hard]} Explain Sandboxing from first principles, then show how you would evaluate and operate it in production.
\noindent\textbf{Answer.} Start with the mechanism: Defense in depth combines process, filesystem, network, syscall, credential, resource, and time boundaries. Make the tradeoff explicit: Strong isolation raises operational cost and may restrict legitimate tools. Define workload-specific offline metrics and a latency/cost budget, test important slices, instrument the critical path, and create a rollback criterion. The production review must explicitly guard against this failure: Running as non-root inside an otherwise privileged container is insufficient.
\par\textit{A strong answer connects mechanism, measurable tradeoffs, failure isolation, observability, and rollback rather than listing product features.}
\paragraph{MCQ-1761 [medium]} Which statement best describes Data privacy?
\begin{enumerate}[label=\Alph*.]
\item Structured schema, semantic checks, policy filters, allowlists, and deterministic verification gate model output.
\item Weights, datasets, containers, packages, adapters, and prompts require provenance, scanning, signing, and review.
\item Data minimization, purpose limitation, access control, retention, deletion, and regional handling apply to prompts and traces.
\item The host checks identity, scope, resource, action, policy, and user intent before every tool execution.
\end{enumerate}
\noindent\textbf{Answer.} Data minimization, purpose limitation, access control, retention, deletion, and regional handling apply to prompts and traces.
\par\textit{Data minimization, purpose limitation, access control, retention, deletion, and regional handling apply to prompts and traces. Tradeoff: Privacy protection can reduce debugging data and personalization. Watch for: Redacting only user messages while retaining retrieved documents leaves exposure.}
\paragraph{MCQ-1762 [hard]} What is the central engineering tradeoff of Data privacy?
\begin{enumerate}[label=\Alph*.]
\item Privacy protection can reduce debugging data and personalization.
\item Assurance improves while promotion pipelines become more deliberate.
\item Fail-closed validation improves safety but may lower completion rate.
\item Fine-grained authorization limits blast radius but adds policy complexity.
\end{enumerate}
\noindent\textbf{Answer.} Privacy protection can reduce debugging data and personalization.
\par\textit{Data minimization, purpose limitation, access control, retention, deletion, and regional handling apply to prompts and traces. Tradeoff: Privacy protection can reduce debugging data and personalization. Watch for: Redacting only user messages while retaining retrieved documents leaves exposure.}
\paragraph{MCQ-1763 [hard]} Which failure mode should an interviewer expect you to identify for Data privacy?
\begin{enumerate}[label=\Alph*.]
\item Loading arbitrary pickle-based artifacts can execute code.
\item A valid JSON schema does not prove safe SQL, URLs, or business intent.
\item Authorizing the agent once at session start enables confused-deputy attacks.
\item Redacting only user messages while retaining retrieved documents leaves exposure.
\end{enumerate}
\noindent\textbf{Answer.} Redacting only user messages while retaining retrieved documents leaves exposure.
\par\textit{Data minimization, purpose limitation, access control, retention, deletion, and regional handling apply to prompts and traces. Tradeoff: Privacy protection can reduce debugging data and personalization. Watch for: Redacting only user messages while retaining retrieved documents leaves exposure.}
\paragraph{MCQ-1764 [medium]} A design review is specifically concerned with this risk: Redacting only user messages while retaining retrieved documents leaves exposure. Which topic should be reviewed first?
\begin{enumerate}[label=\Alph*.]
\item Data privacy
\item Output validation
\item Model supply chain
\item Tool authorization
\end{enumerate}
\noindent\textbf{Answer.} Data privacy
\par\textit{Data minimization, purpose limitation, access control, retention, deletion, and regional handling apply to prompts and traces. Tradeoff: Privacy protection can reduce debugging data and personalization. Watch for: Redacting only user messages while retaining retrieved documents leaves exposure.}
\paragraph{MCQ-1765 [hard]} Which design note most accurately balances the benefits and costs of Data privacy?
\begin{enumerate}[label=\Alph*.]
\item Privacy protection can reduce debugging data and personalization.
\item Fine-grained authorization limits blast radius but adds policy complexity.
\item Fail-closed validation improves safety but may lower completion rate.
\item Assurance improves while promotion pipelines become more deliberate.
\end{enumerate}
\noindent\textbf{Answer.} Privacy protection can reduce debugging data and personalization.
\par\textit{Data minimization, purpose limitation, access control, retention, deletion, and regional handling apply to prompts and traces. Tradeoff: Privacy protection can reduce debugging data and personalization. Watch for: Redacting only user messages while retaining retrieved documents leaves exposure.}
\paragraph{MCQ-1766 [medium]} Which explanation would be strongest in a system-design interview about Data privacy?
\begin{enumerate}[label=\Alph*.]
\item Data minimization, purpose limitation, access control, retention, deletion, and regional handling apply to prompts and traces.
\item The host checks identity, scope, resource, action, policy, and user intent before every tool execution.
\item Weights, datasets, containers, packages, adapters, and prompts require provenance, scanning, signing, and review.
\item Structured schema, semantic checks, policy filters, allowlists, and deterministic verification gate model output.
\end{enumerate}
\noindent\textbf{Answer.} Data minimization, purpose limitation, access control, retention, deletion, and regional handling apply to prompts and traces.
\par\textit{Data minimization, purpose limitation, access control, retention, deletion, and regional handling apply to prompts and traces. Tradeoff: Privacy protection can reduce debugging data and personalization. Watch for: Redacting only user messages while retaining retrieved documents leaves exposure.}
\paragraph{MCQ-1767 [hard]} Which production warning is most directly associated with Data privacy?
\begin{enumerate}[label=\Alph*.]
\item Loading arbitrary pickle-based artifacts can execute code.
\item A valid JSON schema does not prove safe SQL, URLs, or business intent.
\item Authorizing the agent once at session start enables confused-deputy attacks.
\item Redacting only user messages while retaining retrieved documents leaves exposure.
\end{enumerate}
\noindent\textbf{Answer.} Redacting only user messages while retaining retrieved documents leaves exposure.
\par\textit{Data minimization, purpose limitation, access control, retention, deletion, and regional handling apply to prompts and traces. Tradeoff: Privacy protection can reduce debugging data and personalization. Watch for: Redacting only user messages while retaining retrieved documents leaves exposure.}
\paragraph{FC-1768 [medium]} Define Data privacy in interview-ready language.
\noindent\textbf{Answer.} Data minimization, purpose limitation, access control, retention, deletion, and regional handling apply to prompts and traces.
\par\textit{Data minimization, purpose limitation, access control, retention, deletion, and regional handling apply to prompts and traces.}
\paragraph{FC-1769 [hard]} State the main tradeoff and production pitfall for Data privacy.
\noindent\textbf{Answer.} Tradeoff: Privacy protection can reduce debugging data and personalization. Pitfall: Redacting only user messages while retaining retrieved documents leaves exposure.
\par\textit{Tradeoff: Privacy protection can reduce debugging data and personalization. Pitfall: Redacting only user messages while retaining retrieved documents leaves exposure.}
\paragraph{LA-1770 [hard]} Explain Data privacy from first principles, then show how you would evaluate and operate it in production.
\noindent\textbf{Answer.} Start with the mechanism: Data minimization, purpose limitation, access control, retention, deletion, and regional handling apply to prompts and traces. Make the tradeoff explicit: Privacy protection can reduce debugging data and personalization. Define workload-specific offline metrics and a latency/cost budget, test important slices, instrument the critical path, and create a rollback criterion. The production review must explicitly guard against this failure: Redacting only user messages while retaining retrieved documents leaves exposure.
\par\textit{A strong answer connects mechanism, measurable tradeoffs, failure isolation, observability, and rollback rather than listing product features.}
\paragraph{MCQ-1771 [medium]} Which statement best describes Output validation?
\begin{enumerate}[label=\Alph*.]
\item Defense in depth combines process, filesystem, network, syscall, credential, resource, and time boundaries.
\item Weights, datasets, containers, packages, adapters, and prompts require provenance, scanning, signing, and review.
\item Structured schema, semantic checks, policy filters, allowlists, and deterministic verification gate model output.
\item Adversarial testing explores misuse, prompt injection, privacy, bias, jailbreaks, tool abuse, and failure recovery.
\end{enumerate}
\noindent\textbf{Answer.} Structured schema, semantic checks, policy filters, allowlists, and deterministic verification gate model output.
\par\textit{Structured schema, semantic checks, policy filters, allowlists, and deterministic verification gate model output. Tradeoff: Fail-closed validation improves safety but may lower completion rate. Watch for: A valid JSON schema does not prove safe SQL, URLs, or business intent.}
\paragraph{MCQ-1772 [hard]} What is the central engineering tradeoff of Output validation?
\begin{enumerate}[label=\Alph*.]
\item Strong isolation raises operational cost and may restrict legitimate tools.
\item It reveals unknown weaknesses but provides samples, not a guarantee of safety.
\item Fail-closed validation improves safety but may lower completion rate.
\item Assurance improves while promotion pipelines become more deliberate.
\end{enumerate}
\noindent\textbf{Answer.} Fail-closed validation improves safety but may lower completion rate.
\par\textit{Structured schema, semantic checks, policy filters, allowlists, and deterministic verification gate model output. Tradeoff: Fail-closed validation improves safety but may lower completion rate. Watch for: A valid JSON schema does not prove safe SQL, URLs, or business intent.}
\paragraph{MCQ-1773 [hard]} Which failure mode should an interviewer expect you to identify for Output validation?
\begin{enumerate}[label=\Alph*.]
\item A valid JSON schema does not prove safe SQL, URLs, or business intent.
\item Running one generic jailbreak list misses application-specific assets and actions.
\item Loading arbitrary pickle-based artifacts can execute code.
\item Running as non-root inside an otherwise privileged container is insufficient.
\end{enumerate}
\noindent\textbf{Answer.} A valid JSON schema does not prove safe SQL, URLs, or business intent.
\par\textit{Structured schema, semantic checks, policy filters, allowlists, and deterministic verification gate model output. Tradeoff: Fail-closed validation improves safety but may lower completion rate. Watch for: A valid JSON schema does not prove safe SQL, URLs, or business intent.}
\paragraph{MCQ-1774 [medium]} A design review is specifically concerned with this risk: A valid JSON schema does not prove safe SQL, URLs, or business intent. Which topic should be reviewed first?
\begin{enumerate}[label=\Alph*.]
\item Model supply chain
\item Red teaming
\item Sandboxing
\item Output validation
\end{enumerate}
\noindent\textbf{Answer.} Output validation
\par\textit{Structured schema, semantic checks, policy filters, allowlists, and deterministic verification gate model output. Tradeoff: Fail-closed validation improves safety but may lower completion rate. Watch for: A valid JSON schema does not prove safe SQL, URLs, or business intent.}
\paragraph{MCQ-1775 [hard]} Which design note most accurately balances the benefits and costs of Output validation?
\begin{enumerate}[label=\Alph*.]
\item Strong isolation raises operational cost and may restrict legitimate tools.
\item Fail-closed validation improves safety but may lower completion rate.
\item Assurance improves while promotion pipelines become more deliberate.
\item It reveals unknown weaknesses but provides samples, not a guarantee of safety.
\end{enumerate}
\noindent\textbf{Answer.} Fail-closed validation improves safety but may lower completion rate.
\par\textit{Structured schema, semantic checks, policy filters, allowlists, and deterministic verification gate model output. Tradeoff: Fail-closed validation improves safety but may lower completion rate. Watch for: A valid JSON schema does not prove safe SQL, URLs, or business intent.}
\paragraph{MCQ-1776 [medium]} Which explanation would be strongest in a system-design interview about Output validation?
\begin{enumerate}[label=\Alph*.]
\item Weights, datasets, containers, packages, adapters, and prompts require provenance, scanning, signing, and review.
\item Structured schema, semantic checks, policy filters, allowlists, and deterministic verification gate model output.
\item Adversarial testing explores misuse, prompt injection, privacy, bias, jailbreaks, tool abuse, and failure recovery.
\item Defense in depth combines process, filesystem, network, syscall, credential, resource, and time boundaries.
\end{enumerate}
\noindent\textbf{Answer.} Structured schema, semantic checks, policy filters, allowlists, and deterministic verification gate model output.
\par\textit{Structured schema, semantic checks, policy filters, allowlists, and deterministic verification gate model output. Tradeoff: Fail-closed validation improves safety but may lower completion rate. Watch for: A valid JSON schema does not prove safe SQL, URLs, or business intent.}
\paragraph{MCQ-1777 [hard]} Which production warning is most directly associated with Output validation?
\begin{enumerate}[label=\Alph*.]
\item Loading arbitrary pickle-based artifacts can execute code.
\item A valid JSON schema does not prove safe SQL, URLs, or business intent.
\item Running one generic jailbreak list misses application-specific assets and actions.
\item Running as non-root inside an otherwise privileged container is insufficient.
\end{enumerate}
\noindent\textbf{Answer.} A valid JSON schema does not prove safe SQL, URLs, or business intent.
\par\textit{Structured schema, semantic checks, policy filters, allowlists, and deterministic verification gate model output. Tradeoff: Fail-closed validation improves safety but may lower completion rate. Watch for: A valid JSON schema does not prove safe SQL, URLs, or business intent.}
\paragraph{FC-1778 [medium]} Define Output validation in interview-ready language.
\noindent\textbf{Answer.} Structured schema, semantic checks, policy filters, allowlists, and deterministic verification gate model output.
\par\textit{Structured schema, semantic checks, policy filters, allowlists, and deterministic verification gate model output.}
\paragraph{FC-1779 [hard]} State the main tradeoff and production pitfall for Output validation.
\noindent\textbf{Answer.} Tradeoff: Fail-closed validation improves safety but may lower completion rate. Pitfall: A valid JSON schema does not prove safe SQL, URLs, or business intent.
\par\textit{Tradeoff: Fail-closed validation improves safety but may lower completion rate. Pitfall: A valid JSON schema does not prove safe SQL, URLs, or business intent.}
\paragraph{LA-1780 [hard]} Explain Output validation from first principles, then show how you would evaluate and operate it in production.
\noindent\textbf{Answer.} Start with the mechanism: Structured schema, semantic checks, policy filters, allowlists, and deterministic verification gate model output. Make the tradeoff explicit: Fail-closed validation improves safety but may lower completion rate. Define workload-specific offline metrics and a latency/cost budget, test important slices, instrument the critical path, and create a rollback criterion. The production review must explicitly guard against this failure: A valid JSON schema does not prove safe SQL, URLs, or business intent.
\par\textit{A strong answer connects mechanism, measurable tradeoffs, failure isolation, observability, and rollback rather than listing product features.}
\paragraph{MCQ-1781 [medium]} Which statement best describes Model supply chain?
\begin{enumerate}[label=\Alph*.]
\item The host checks identity, scope, resource, action, policy, and user intent before every tool execution.
\item Risk processes map context, measure impacts, manage controls, document evidence, and govern residual risk.
\item Credentials are issued narrowly, stored outside prompts and code, rotated, audited, and redacted from telemetry.
\item Weights, datasets, containers, packages, adapters, and prompts require provenance, scanning, signing, and review.
\end{enumerate}
\noindent\textbf{Answer.} Weights, datasets, containers, packages, adapters, and prompts require provenance, scanning, signing, and review.
\par\textit{Weights, datasets, containers, packages, adapters, and prompts require provenance, scanning, signing, and review. Tradeoff: Assurance improves while promotion pipelines become more deliberate. Watch for: Loading arbitrary pickle-based artifacts can execute code.}
\paragraph{MCQ-1782 [hard]} What is the central engineering tradeoff of Model supply chain?
\begin{enumerate}[label=\Alph*.]
\item Governance improves accountability but can become checkbox compliance without technical signals.
\item Assurance improves while promotion pipelines become more deliberate.
\item Short-lived credentials reduce exposure but require reliable identity infrastructure.
\item Fine-grained authorization limits blast radius but adds policy complexity.
\end{enumerate}
\noindent\textbf{Answer.} Assurance improves while promotion pipelines become more deliberate.
\par\textit{Weights, datasets, containers, packages, adapters, and prompts require provenance, scanning, signing, and review. Tradeoff: Assurance improves while promotion pipelines become more deliberate. Watch for: Loading arbitrary pickle-based artifacts can execute code.}
\paragraph{MCQ-1783 [hard]} Which failure mode should an interviewer expect you to identify for Model supply chain?
\begin{enumerate}[label=\Alph*.]
\item Authorizing the agent once at session start enables confused-deputy attacks.
\item Environment variables copied into sandboxes or crash logs leak broad authority.
\item Loading arbitrary pickle-based artifacts can execute code.
\item A model card without deployment-specific threat modeling is incomplete.
\end{enumerate}
\noindent\textbf{Answer.} Loading arbitrary pickle-based artifacts can execute code.
\par\textit{Weights, datasets, containers, packages, adapters, and prompts require provenance, scanning, signing, and review. Tradeoff: Assurance improves while promotion pipelines become more deliberate. Watch for: Loading arbitrary pickle-based artifacts can execute code.}
\paragraph{MCQ-1784 [medium]} A design review is specifically concerned with this risk: Loading arbitrary pickle-based artifacts can execute code. Which topic should be reviewed first?
\begin{enumerate}[label=\Alph*.]
\item Tool authorization
\item Model supply chain
\item Secrets management
\item AI risk management
\end{enumerate}
\noindent\textbf{Answer.} Model supply chain
\par\textit{Weights, datasets, containers, packages, adapters, and prompts require provenance, scanning, signing, and review. Tradeoff: Assurance improves while promotion pipelines become more deliberate. Watch for: Loading arbitrary pickle-based artifacts can execute code.}
\paragraph{MCQ-1785 [hard]} Which design note most accurately balances the benefits and costs of Model supply chain?
\begin{enumerate}[label=\Alph*.]
\item Short-lived credentials reduce exposure but require reliable identity infrastructure.
\item Fine-grained authorization limits blast radius but adds policy complexity.
\item Assurance improves while promotion pipelines become more deliberate.
\item Governance improves accountability but can become checkbox compliance without technical signals.
\end{enumerate}
\noindent\textbf{Answer.} Assurance improves while promotion pipelines become more deliberate.
\par\textit{Weights, datasets, containers, packages, adapters, and prompts require provenance, scanning, signing, and review. Tradeoff: Assurance improves while promotion pipelines become more deliberate. Watch for: Loading arbitrary pickle-based artifacts can execute code.}
\paragraph{MCQ-1786 [medium]} Which explanation would be strongest in a system-design interview about Model supply chain?
\begin{enumerate}[label=\Alph*.]
\item Risk processes map context, measure impacts, manage controls, document evidence, and govern residual risk.
\item The host checks identity, scope, resource, action, policy, and user intent before every tool execution.
\item Weights, datasets, containers, packages, adapters, and prompts require provenance, scanning, signing, and review.
\item Credentials are issued narrowly, stored outside prompts and code, rotated, audited, and redacted from telemetry.
\end{enumerate}
\noindent\textbf{Answer.} Weights, datasets, containers, packages, adapters, and prompts require provenance, scanning, signing, and review.
\par\textit{Weights, datasets, containers, packages, adapters, and prompts require provenance, scanning, signing, and review. Tradeoff: Assurance improves while promotion pipelines become more deliberate. Watch for: Loading arbitrary pickle-based artifacts can execute code.}
\paragraph{MCQ-1787 [hard]} Which production warning is most directly associated with Model supply chain?
\begin{enumerate}[label=\Alph*.]
\item A model card without deployment-specific threat modeling is incomplete.
\item Environment variables copied into sandboxes or crash logs leak broad authority.
\item Authorizing the agent once at session start enables confused-deputy attacks.
\item Loading arbitrary pickle-based artifacts can execute code.
\end{enumerate}
\noindent\textbf{Answer.} Loading arbitrary pickle-based artifacts can execute code.
\par\textit{Weights, datasets, containers, packages, adapters, and prompts require provenance, scanning, signing, and review. Tradeoff: Assurance improves while promotion pipelines become more deliberate. Watch for: Loading arbitrary pickle-based artifacts can execute code.}
\paragraph{FC-1788 [medium]} Define Model supply chain in interview-ready language.
\noindent\textbf{Answer.} Weights, datasets, containers, packages, adapters, and prompts require provenance, scanning, signing, and review.
\par\textit{Weights, datasets, containers, packages, adapters, and prompts require provenance, scanning, signing, and review.}
\paragraph{FC-1789 [hard]} State the main tradeoff and production pitfall for Model supply chain.
\noindent\textbf{Answer.} Tradeoff: Assurance improves while promotion pipelines become more deliberate. Pitfall: Loading arbitrary pickle-based artifacts can execute code.
\par\textit{Tradeoff: Assurance improves while promotion pipelines become more deliberate. Pitfall: Loading arbitrary pickle-based artifacts can execute code.}
\paragraph{LA-1790 [hard]} Explain Model supply chain from first principles, then show how you would evaluate and operate it in production.
\noindent\textbf{Answer.} Start with the mechanism: Weights, datasets, containers, packages, adapters, and prompts require provenance, scanning, signing, and review. Make the tradeoff explicit: Assurance improves while promotion pipelines become more deliberate. Define workload-specific offline metrics and a latency/cost budget, test important slices, instrument the critical path, and create a rollback criterion. The production review must explicitly guard against this failure: Loading arbitrary pickle-based artifacts can execute code.
\par\textit{A strong answer connects mechanism, measurable tradeoffs, failure isolation, observability, and rollback rather than listing product features.}
\paragraph{MCQ-1791 [medium]} Which statement best describes AI risk management?
\begin{enumerate}[label=\Alph*.]
\item Risk processes map context, measure impacts, manage controls, document evidence, and govern residual risk.
\item The host checks identity, scope, resource, action, policy, and user intent before every tool execution.
\item Immutable event records connect identity, input, policy, model, evidence, tool calls, decisions, and outcomes.
\item Data minimization, purpose limitation, access control, retention, deletion, and regional handling apply to prompts and traces.
\end{enumerate}
\noindent\textbf{Answer.} Risk processes map context, measure impacts, manage controls, document evidence, and govern residual risk.
\par\textit{Risk processes map context, measure impacts, manage controls, document evidence, and govern residual risk. Tradeoff: Governance improves accountability but can become checkbox compliance without technical signals. Watch for: A model card without deployment-specific threat modeling is incomplete.}
\paragraph{MCQ-1792 [hard]} What is the central engineering tradeoff of AI risk management?
\begin{enumerate}[label=\Alph*.]
\item Governance improves accountability but can become checkbox compliance without technical signals.
\item Fine-grained authorization limits blast radius but adds policy complexity.
\item Privacy protection can reduce debugging data and personalization.
\item Investigations and compliance improve, while logs themselves become sensitive assets.
\end{enumerate}
\noindent\textbf{Answer.} Governance improves accountability but can become checkbox compliance without technical signals.
\par\textit{Risk processes map context, measure impacts, manage controls, document evidence, and govern residual risk. Tradeoff: Governance improves accountability but can become checkbox compliance without technical signals. Watch for: A model card without deployment-specific threat modeling is incomplete.}
\paragraph{MCQ-1793 [hard]} Which failure mode should an interviewer expect you to identify for AI risk management?
\begin{enumerate}[label=\Alph*.]
\item Authorizing the agent once at session start enables confused-deputy attacks.
\item A model card without deployment-specific threat modeling is incomplete.
\item Redacting only user messages while retaining retrieved documents leaves exposure.
\item Logging reasoning text instead of decisive actions and evidence can reduce clarity.
\end{enumerate}
\noindent\textbf{Answer.} A model card without deployment-specific threat modeling is incomplete.
\par\textit{Risk processes map context, measure impacts, manage controls, document evidence, and govern residual risk. Tradeoff: Governance improves accountability but can become checkbox compliance without technical signals. Watch for: A model card without deployment-specific threat modeling is incomplete.}
\paragraph{MCQ-1794 [medium]} A design review is specifically concerned with this risk: A model card without deployment-specific threat modeling is incomplete. Which topic should be reviewed first?
\begin{enumerate}[label=\Alph*.]
\item Data privacy
\item AI risk management
\item Tool authorization
\item Auditability
\end{enumerate}
\noindent\textbf{Answer.} AI risk management
\par\textit{Risk processes map context, measure impacts, manage controls, document evidence, and govern residual risk. Tradeoff: Governance improves accountability but can become checkbox compliance without technical signals. Watch for: A model card without deployment-specific threat modeling is incomplete.}
\paragraph{MCQ-1795 [hard]} Which design note most accurately balances the benefits and costs of AI risk management?
\begin{enumerate}[label=\Alph*.]
\item Investigations and compliance improve, while logs themselves become sensitive assets.
\item Privacy protection can reduce debugging data and personalization.
\item Fine-grained authorization limits blast radius but adds policy complexity.
\item Governance improves accountability but can become checkbox compliance without technical signals.
\end{enumerate}
\noindent\textbf{Answer.} Governance improves accountability but can become checkbox compliance without technical signals.
\par\textit{Risk processes map context, measure impacts, manage controls, document evidence, and govern residual risk. Tradeoff: Governance improves accountability but can become checkbox compliance without technical signals. Watch for: A model card without deployment-specific threat modeling is incomplete.}
\paragraph{MCQ-1796 [medium]} Which explanation would be strongest in a system-design interview about AI risk management?
\begin{enumerate}[label=\Alph*.]
\item Immutable event records connect identity, input, policy, model, evidence, tool calls, decisions, and outcomes.
\item Data minimization, purpose limitation, access control, retention, deletion, and regional handling apply to prompts and traces.
\item The host checks identity, scope, resource, action, policy, and user intent before every tool execution.
\item Risk processes map context, measure impacts, manage controls, document evidence, and govern residual risk.
\end{enumerate}
\noindent\textbf{Answer.} Risk processes map context, measure impacts, manage controls, document evidence, and govern residual risk.
\par\textit{Risk processes map context, measure impacts, manage controls, document evidence, and govern residual risk. Tradeoff: Governance improves accountability but can become checkbox compliance without technical signals. Watch for: A model card without deployment-specific threat modeling is incomplete.}
\paragraph{MCQ-1797 [hard]} Which production warning is most directly associated with AI risk management?
\begin{enumerate}[label=\Alph*.]
\item Logging reasoning text instead of decisive actions and evidence can reduce clarity.
\item A model card without deployment-specific threat modeling is incomplete.
\item Redacting only user messages while retaining retrieved documents leaves exposure.
\item Authorizing the agent once at session start enables confused-deputy attacks.
\end{enumerate}
\noindent\textbf{Answer.} A model card without deployment-specific threat modeling is incomplete.
\par\textit{Risk processes map context, measure impacts, manage controls, document evidence, and govern residual risk. Tradeoff: Governance improves accountability but can become checkbox compliance without technical signals. Watch for: A model card without deployment-specific threat modeling is incomplete.}
\paragraph{FC-1798 [medium]} Define AI risk management in interview-ready language.
\noindent\textbf{Answer.} Risk processes map context, measure impacts, manage controls, document evidence, and govern residual risk.
\par\textit{Risk processes map context, measure impacts, manage controls, document evidence, and govern residual risk.}
\paragraph{FC-1799 [hard]} State the main tradeoff and production pitfall for AI risk management.
\noindent\textbf{Answer.} Tradeoff: Governance improves accountability but can become checkbox compliance without technical signals. Pitfall: A model card without deployment-specific threat modeling is incomplete.
\par\textit{Tradeoff: Governance improves accountability but can become checkbox compliance without technical signals. Pitfall: A model card without deployment-specific threat modeling is incomplete.}
\paragraph{LA-1800 [hard]} Explain AI risk management from first principles, then show how you would evaluate and operate it in production.
\noindent\textbf{Answer.} Start with the mechanism: Risk processes map context, measure impacts, manage controls, document evidence, and govern residual risk. Make the tradeoff explicit: Governance improves accountability but can become checkbox compliance without technical signals. Define workload-specific offline metrics and a latency/cost budget, test important slices, instrument the critical path, and create a rollback criterion. The production review must explicitly guard against this failure: A model card without deployment-specific threat modeling is incomplete.
\par\textit{A strong answer connects mechanism, measurable tradeoffs, failure isolation, observability, and rollback rather than listing product features.}
\paragraph{MCQ-1801 [medium]} Which statement best describes Red teaming?
\begin{enumerate}[label=\Alph*.]
\item Risk processes map context, measure impacts, manage controls, document evidence, and govern residual risk.
\item Immutable event records connect identity, input, policy, model, evidence, tool calls, decisions, and outcomes.
\item Data minimization, purpose limitation, access control, retention, deletion, and regional handling apply to prompts and traces.
\item Adversarial testing explores misuse, prompt injection, privacy, bias, jailbreaks, tool abuse, and failure recovery.
\end{enumerate}
\noindent\textbf{Answer.} Adversarial testing explores misuse, prompt injection, privacy, bias, jailbreaks, tool abuse, and failure recovery.
\par\textit{Adversarial testing explores misuse, prompt injection, privacy, bias, jailbreaks, tool abuse, and failure recovery. Tradeoff: It reveals unknown weaknesses but provides samples, not a guarantee of safety. Watch for: Running one generic jailbreak list misses application-specific assets and actions.}
\paragraph{MCQ-1802 [hard]} What is the central engineering tradeoff of Red teaming?
\begin{enumerate}[label=\Alph*.]
\item Investigations and compliance improve, while logs themselves become sensitive assets.
\item It reveals unknown weaknesses but provides samples, not a guarantee of safety.
\item Privacy protection can reduce debugging data and personalization.
\item Governance improves accountability but can become checkbox compliance without technical signals.
\end{enumerate}
\noindent\textbf{Answer.} It reveals unknown weaknesses but provides samples, not a guarantee of safety.
\par\textit{Adversarial testing explores misuse, prompt injection, privacy, bias, jailbreaks, tool abuse, and failure recovery. Tradeoff: It reveals unknown weaknesses but provides samples, not a guarantee of safety. Watch for: Running one generic jailbreak list misses application-specific assets and actions.}
\paragraph{MCQ-1803 [hard]} Which failure mode should an interviewer expect you to identify for Red teaming?
\begin{enumerate}[label=\Alph*.]
\item Redacting only user messages while retaining retrieved documents leaves exposure.
\item Logging reasoning text instead of decisive actions and evidence can reduce clarity.
\item Running one generic jailbreak list misses application-specific assets and actions.
\item A model card without deployment-specific threat modeling is incomplete.
\end{enumerate}
\noindent\textbf{Answer.} Running one generic jailbreak list misses application-specific assets and actions.
\par\textit{Adversarial testing explores misuse, prompt injection, privacy, bias, jailbreaks, tool abuse, and failure recovery. Tradeoff: It reveals unknown weaknesses but provides samples, not a guarantee of safety. Watch for: Running one generic jailbreak list misses application-specific assets and actions.}
\paragraph{MCQ-1804 [medium]} A design review is specifically concerned with this risk: Running one generic jailbreak list misses application-specific assets and actions. Which topic should be reviewed first?
\begin{enumerate}[label=\Alph*.]
\item Red teaming
\item AI risk management
\item Auditability
\item Data privacy
\end{enumerate}
\noindent\textbf{Answer.} Red teaming
\par\textit{Adversarial testing explores misuse, prompt injection, privacy, bias, jailbreaks, tool abuse, and failure recovery. Tradeoff: It reveals unknown weaknesses but provides samples, not a guarantee of safety. Watch for: Running one generic jailbreak list misses application-specific assets and actions.}
\paragraph{MCQ-1805 [hard]} Which design note most accurately balances the benefits and costs of Red teaming?
\begin{enumerate}[label=\Alph*.]
\item Governance improves accountability but can become checkbox compliance without technical signals.
\item It reveals unknown weaknesses but provides samples, not a guarantee of safety.
\item Privacy protection can reduce debugging data and personalization.
\item Investigations and compliance improve, while logs themselves become sensitive assets.
\end{enumerate}
\noindent\textbf{Answer.} It reveals unknown weaknesses but provides samples, not a guarantee of safety.
\par\textit{Adversarial testing explores misuse, prompt injection, privacy, bias, jailbreaks, tool abuse, and failure recovery. Tradeoff: It reveals unknown weaknesses but provides samples, not a guarantee of safety. Watch for: Running one generic jailbreak list misses application-specific assets and actions.}
\paragraph{MCQ-1806 [medium]} Which explanation would be strongest in a system-design interview about Red teaming?
\begin{enumerate}[label=\Alph*.]
\item Adversarial testing explores misuse, prompt injection, privacy, bias, jailbreaks, tool abuse, and failure recovery.
\item Risk processes map context, measure impacts, manage controls, document evidence, and govern residual risk.
\item Data minimization, purpose limitation, access control, retention, deletion, and regional handling apply to prompts and traces.
\item Immutable event records connect identity, input, policy, model, evidence, tool calls, decisions, and outcomes.
\end{enumerate}
\noindent\textbf{Answer.} Adversarial testing explores misuse, prompt injection, privacy, bias, jailbreaks, tool abuse, and failure recovery.
\par\textit{Adversarial testing explores misuse, prompt injection, privacy, bias, jailbreaks, tool abuse, and failure recovery. Tradeoff: It reveals unknown weaknesses but provides samples, not a guarantee of safety. Watch for: Running one generic jailbreak list misses application-specific assets and actions.}
\paragraph{MCQ-1807 [hard]} Which production warning is most directly associated with Red teaming?
\begin{enumerate}[label=\Alph*.]
\item A model card without deployment-specific threat modeling is incomplete.
\item Redacting only user messages while retaining retrieved documents leaves exposure.
\item Logging reasoning text instead of decisive actions and evidence can reduce clarity.
\item Running one generic jailbreak list misses application-specific assets and actions.
\end{enumerate}
\noindent\textbf{Answer.} Running one generic jailbreak list misses application-specific assets and actions.
\par\textit{Adversarial testing explores misuse, prompt injection, privacy, bias, jailbreaks, tool abuse, and failure recovery. Tradeoff: It reveals unknown weaknesses but provides samples, not a guarantee of safety. Watch for: Running one generic jailbreak list misses application-specific assets and actions.}
\paragraph{FC-1808 [medium]} Define Red teaming in interview-ready language.
\noindent\textbf{Answer.} Adversarial testing explores misuse, prompt injection, privacy, bias, jailbreaks, tool abuse, and failure recovery.
\par\textit{Adversarial testing explores misuse, prompt injection, privacy, bias, jailbreaks, tool abuse, and failure recovery.}
\paragraph{FC-1809 [hard]} State the main tradeoff and production pitfall for Red teaming.
\noindent\textbf{Answer.} Tradeoff: It reveals unknown weaknesses but provides samples, not a guarantee of safety. Pitfall: Running one generic jailbreak list misses application-specific assets and actions.
\par\textit{Tradeoff: It reveals unknown weaknesses but provides samples, not a guarantee of safety. Pitfall: Running one generic jailbreak list misses application-specific assets and actions.}
\paragraph{LA-1810 [hard]} Explain Red teaming from first principles, then show how you would evaluate and operate it in production.
\noindent\textbf{Answer.} Start with the mechanism: Adversarial testing explores misuse, prompt injection, privacy, bias, jailbreaks, tool abuse, and failure recovery. Make the tradeoff explicit: It reveals unknown weaknesses but provides samples, not a guarantee of safety. Define workload-specific offline metrics and a latency/cost budget, test important slices, instrument the critical path, and create a rollback criterion. The production review must explicitly guard against this failure: Running one generic jailbreak list misses application-specific assets and actions.
\par\textit{A strong answer connects mechanism, measurable tradeoffs, failure isolation, observability, and rollback rather than listing product features.}
\paragraph{MCQ-1811 [medium]} Which statement best describes Auditability?
\begin{enumerate}[label=\Alph*.]
\item The host checks identity, scope, resource, action, policy, and user intent before every tool execution.
\item Immutable event records connect identity, input, policy, model, evidence, tool calls, decisions, and outcomes.
\item Untrusted content attempts to override instructions, exfiltrate data, or induce tool actions.
\item Risk processes map context, measure impacts, manage controls, document evidence, and govern residual risk.
\end{enumerate}
\noindent\textbf{Answer.} Immutable event records connect identity, input, policy, model, evidence, tool calls, decisions, and outcomes.
\par\textit{Immutable event records connect identity, input, policy, model, evidence, tool calls, decisions, and outcomes. Tradeoff: Investigations and compliance improve, while logs themselves become sensitive assets. Watch for: Logging reasoning text instead of decisive actions and evidence can reduce clarity.}
\paragraph{MCQ-1812 [hard]} What is the central engineering tradeoff of Auditability?
\begin{enumerate}[label=\Alph*.]
\item Governance improves accountability but can become checkbox compliance without technical signals.
\item Fine-grained authorization limits blast radius but adds policy complexity.
\item Investigations and compliance improve, while logs themselves become sensitive assets.
\item Content isolation and least privilege reduce impact but cannot make untrusted text authoritative.
\end{enumerate}
\noindent\textbf{Answer.} Investigations and compliance improve, while logs themselves become sensitive assets.
\par\textit{Immutable event records connect identity, input, policy, model, evidence, tool calls, decisions, and outcomes. Tradeoff: Investigations and compliance improve, while logs themselves become sensitive assets. Watch for: Logging reasoning text instead of decisive actions and evidence can reduce clarity.}
\paragraph{MCQ-1813 [hard]} Which failure mode should an interviewer expect you to identify for Auditability?
\begin{enumerate}[label=\Alph*.]
\item Telling the model to ignore attacks is not a security boundary.
\item Logging reasoning text instead of decisive actions and evidence can reduce clarity.
\item A model card without deployment-specific threat modeling is incomplete.
\item Authorizing the agent once at session start enables confused-deputy attacks.
\end{enumerate}
\noindent\textbf{Answer.} Logging reasoning text instead of decisive actions and evidence can reduce clarity.
\par\textit{Immutable event records connect identity, input, policy, model, evidence, tool calls, decisions, and outcomes. Tradeoff: Investigations and compliance improve, while logs themselves become sensitive assets. Watch for: Logging reasoning text instead of decisive actions and evidence can reduce clarity.}
\paragraph{MCQ-1814 [medium]} A design review is specifically concerned with this risk: Logging reasoning text instead of decisive actions and evidence can reduce clarity. Which topic should be reviewed first?
\begin{enumerate}[label=\Alph*.]
\item AI risk management
\item Prompt injection
\item Auditability
\item Tool authorization
\end{enumerate}
\noindent\textbf{Answer.} Auditability
\par\textit{Immutable event records connect identity, input, policy, model, evidence, tool calls, decisions, and outcomes. Tradeoff: Investigations and compliance improve, while logs themselves become sensitive assets. Watch for: Logging reasoning text instead of decisive actions and evidence can reduce clarity.}
\paragraph{MCQ-1815 [hard]} Which design note most accurately balances the benefits and costs of Auditability?
\begin{enumerate}[label=\Alph*.]
\item Governance improves accountability but can become checkbox compliance without technical signals.
\item Fine-grained authorization limits blast radius but adds policy complexity.
\item Investigations and compliance improve, while logs themselves become sensitive assets.
\item Content isolation and least privilege reduce impact but cannot make untrusted text authoritative.
\end{enumerate}
\noindent\textbf{Answer.} Investigations and compliance improve, while logs themselves become sensitive assets.
\par\textit{Immutable event records connect identity, input, policy, model, evidence, tool calls, decisions, and outcomes. Tradeoff: Investigations and compliance improve, while logs themselves become sensitive assets. Watch for: Logging reasoning text instead of decisive actions and evidence can reduce clarity.}
\paragraph{MCQ-1816 [medium]} Which explanation would be strongest in a system-design interview about Auditability?
\begin{enumerate}[label=\Alph*.]
\item The host checks identity, scope, resource, action, policy, and user intent before every tool execution.
\item Immutable event records connect identity, input, policy, model, evidence, tool calls, decisions, and outcomes.
\item Untrusted content attempts to override instructions, exfiltrate data, or induce tool actions.
\item Risk processes map context, measure impacts, manage controls, document evidence, and govern residual risk.
\end{enumerate}
\noindent\textbf{Answer.} Immutable event records connect identity, input, policy, model, evidence, tool calls, decisions, and outcomes.
\par\textit{Immutable event records connect identity, input, policy, model, evidence, tool calls, decisions, and outcomes. Tradeoff: Investigations and compliance improve, while logs themselves become sensitive assets. Watch for: Logging reasoning text instead of decisive actions and evidence can reduce clarity.}
\paragraph{MCQ-1817 [hard]} Which production warning is most directly associated with Auditability?
\begin{enumerate}[label=\Alph*.]
\item A model card without deployment-specific threat modeling is incomplete.
\item Logging reasoning text instead of decisive actions and evidence can reduce clarity.
\item Telling the model to ignore attacks is not a security boundary.
\item Authorizing the agent once at session start enables confused-deputy attacks.
\end{enumerate}
\noindent\textbf{Answer.} Logging reasoning text instead of decisive actions and evidence can reduce clarity.
\par\textit{Immutable event records connect identity, input, policy, model, evidence, tool calls, decisions, and outcomes. Tradeoff: Investigations and compliance improve, while logs themselves become sensitive assets. Watch for: Logging reasoning text instead of decisive actions and evidence can reduce clarity.}
\paragraph{FC-1818 [medium]} Define Auditability in interview-ready language.
\noindent\textbf{Answer.} Immutable event records connect identity, input, policy, model, evidence, tool calls, decisions, and outcomes.
\par\textit{Immutable event records connect identity, input, policy, model, evidence, tool calls, decisions, and outcomes.}
\paragraph{FC-1819 [hard]} State the main tradeoff and production pitfall for Auditability.
\noindent\textbf{Answer.} Tradeoff: Investigations and compliance improve, while logs themselves become sensitive assets. Pitfall: Logging reasoning text instead of decisive actions and evidence can reduce clarity.
\par\textit{Tradeoff: Investigations and compliance improve, while logs themselves become sensitive assets. Pitfall: Logging reasoning text instead of decisive actions and evidence can reduce clarity.}
\paragraph{LA-1820 [hard]} Explain Auditability from first principles, then show how you would evaluate and operate it in production.
\noindent\textbf{Answer.} Start with the mechanism: Immutable event records connect identity, input, policy, model, evidence, tool calls, decisions, and outcomes. Make the tradeoff explicit: Investigations and compliance improve, while logs themselves become sensitive assets. Define workload-specific offline metrics and a latency/cost budget, test important slices, instrument the critical path, and create a rollback criterion. The production review must explicitly guard against this failure: Logging reasoning text instead of decisive actions and evidence can reduce clarity.
\par\textit{A strong answer connects mechanism, measurable tradeoffs, failure isolation, observability, and rollback rather than listing product features.}
\subsection{Distributed Systems and Reliability}
\paragraph{MCQ-1821 [medium]} Which statement best describes CAP and consistency?
\begin{enumerate}[label=\Alph*.]
\item Bounded queues absorb bursts while admission control and backpressure prevent producers from overwhelming consumers.
\item Balancers route by health, load, locality, affinity, or capacity across replicas.
\item Timeouts bound uncertainty; retries use budgets, backoff, jitter, and idempotency for transient failures.
\item Under a network partition, a distributed system must trade availability against consistent reads and writes.
\end{enumerate}
\noindent\textbf{Answer.} Under a network partition, a distributed system must trade availability against consistent reads and writes.
\par\textit{Under a network partition, a distributed system must trade availability against consistent reads and writes. Tradeoff: Stronger guarantees simplify application invariants but increase latency and unavailability. Watch for: Quoting CAP without identifying partitions and operations does not guide design.}
\paragraph{MCQ-1822 [hard]} What is the central engineering tradeoff of CAP and consistency?
\begin{enumerate}[label=\Alph*.]
\item Success rate improves, while retries amplify overload and tail latency.
\item Stronger guarantees simplify application invariants but increase latency and unavailability.
\item Distribution improves availability but generic least-connections may ignore token workload size.
\item Smoothing improves utilization, while waiting increases latency and stale work.
\end{enumerate}
\noindent\textbf{Answer.} Stronger guarantees simplify application invariants but increase latency and unavailability.
\par\textit{Under a network partition, a distributed system must trade availability against consistent reads and writes. Tradeoff: Stronger guarantees simplify application invariants but increase latency and unavailability. Watch for: Quoting CAP without identifying partitions and operations does not guide design.}
\paragraph{MCQ-1823 [hard]} Which failure mode should an interviewer expect you to identify for CAP and consistency?
\begin{enumerate}[label=\Alph*.]
\item An unbounded queue converts overload into memory exhaustion.
\item Layered automatic retries create multiplicative request storms.
\item Quoting CAP without identifying partitions and operations does not guide design.
\item Sending a long-context request to the least-busy replica can still cause head-of-line blocking.
\end{enumerate}
\noindent\textbf{Answer.} Quoting CAP without identifying partitions and operations does not guide design.
\par\textit{Under a network partition, a distributed system must trade availability against consistent reads and writes. Tradeoff: Stronger guarantees simplify application invariants but increase latency and unavailability. Watch for: Quoting CAP without identifying partitions and operations does not guide design.}
\paragraph{MCQ-1824 [medium]} A design review is specifically concerned with this risk: Quoting CAP without identifying partitions and operations does not guide design. Which topic should be reviewed first?
\begin{enumerate}[label=\Alph*.]
\item Queues and backpressure
\item Load balancing
\item Retries and timeouts
\item CAP and consistency
\end{enumerate}
\noindent\textbf{Answer.} CAP and consistency
\par\textit{Under a network partition, a distributed system must trade availability against consistent reads and writes. Tradeoff: Stronger guarantees simplify application invariants but increase latency and unavailability. Watch for: Quoting CAP without identifying partitions and operations does not guide design.}
\paragraph{MCQ-1825 [hard]} Which design note most accurately balances the benefits and costs of CAP and consistency?
\begin{enumerate}[label=\Alph*.]
\item Stronger guarantees simplify application invariants but increase latency and unavailability.
\item Smoothing improves utilization, while waiting increases latency and stale work.
\item Success rate improves, while retries amplify overload and tail latency.
\item Distribution improves availability but generic least-connections may ignore token workload size.
\end{enumerate}
\noindent\textbf{Answer.} Stronger guarantees simplify application invariants but increase latency and unavailability.
\par\textit{Under a network partition, a distributed system must trade availability against consistent reads and writes. Tradeoff: Stronger guarantees simplify application invariants but increase latency and unavailability. Watch for: Quoting CAP without identifying partitions and operations does not guide design.}
\paragraph{MCQ-1826 [medium]} Which explanation would be strongest in a system-design interview about CAP and consistency?
\begin{enumerate}[label=\Alph*.]
\item Timeouts bound uncertainty; retries use budgets, backoff, jitter, and idempotency for transient failures.
\item Under a network partition, a distributed system must trade availability against consistent reads and writes.
\item Balancers route by health, load, locality, affinity, or capacity across replicas.
\item Bounded queues absorb bursts while admission control and backpressure prevent producers from overwhelming consumers.
\end{enumerate}
\noindent\textbf{Answer.} Under a network partition, a distributed system must trade availability against consistent reads and writes.
\par\textit{Under a network partition, a distributed system must trade availability against consistent reads and writes. Tradeoff: Stronger guarantees simplify application invariants but increase latency and unavailability. Watch for: Quoting CAP without identifying partitions and operations does not guide design.}
\paragraph{MCQ-1827 [hard]} Which production warning is most directly associated with CAP and consistency?
\begin{enumerate}[label=\Alph*.]
\item Quoting CAP without identifying partitions and operations does not guide design.
\item Sending a long-context request to the least-busy replica can still cause head-of-line blocking.
\item Layered automatic retries create multiplicative request storms.
\item An unbounded queue converts overload into memory exhaustion.
\end{enumerate}
\noindent\textbf{Answer.} Quoting CAP without identifying partitions and operations does not guide design.
\par\textit{Under a network partition, a distributed system must trade availability against consistent reads and writes. Tradeoff: Stronger guarantees simplify application invariants but increase latency and unavailability. Watch for: Quoting CAP without identifying partitions and operations does not guide design.}
\paragraph{FC-1828 [medium]} Define CAP and consistency in interview-ready language.
\noindent\textbf{Answer.} Under a network partition, a distributed system must trade availability against consistent reads and writes.
\par\textit{Under a network partition, a distributed system must trade availability against consistent reads and writes.}
\paragraph{FC-1829 [hard]} State the main tradeoff and production pitfall for CAP and consistency.
\noindent\textbf{Answer.} Tradeoff: Stronger guarantees simplify application invariants but increase latency and unavailability. Pitfall: Quoting CAP without identifying partitions and operations does not guide design.
\par\textit{Tradeoff: Stronger guarantees simplify application invariants but increase latency and unavailability. Pitfall: Quoting CAP without identifying partitions and operations does not guide design.}
\paragraph{LA-1830 [hard]} Explain CAP and consistency from first principles, then show how you would evaluate and operate it in production.
\noindent\textbf{Answer.} Start with the mechanism: Under a network partition, a distributed system must trade availability against consistent reads and writes. Make the tradeoff explicit: Stronger guarantees simplify application invariants but increase latency and unavailability. Define workload-specific offline metrics and a latency/cost budget, test important slices, instrument the critical path, and create a rollback criterion. The production review must explicitly guard against this failure: Quoting CAP without identifying partitions and operations does not guide design.
\par\textit{A strong answer connects mechanism, measurable tradeoffs, failure isolation, observability, and rollback rather than listing product features.}
\paragraph{MCQ-1831 [medium]} Which statement best describes Consensus?
\begin{enumerate}[label=\Alph*.]
\item Under a network partition, a distributed system must trade availability against consistent reads and writes.
\item Workload models connect arrival rate, prompt and output tokens, service time, concurrency, memory, and SLO targets.
\item Protocols such as Raft replicate an ordered log so nodes agree despite failures.
\item Timeouts bound uncertainty; retries use budgets, backoff, jitter, and idempotency for transient failures.
\end{enumerate}
\noindent\textbf{Answer.} Protocols such as Raft replicate an ordered log so nodes agree despite failures.
\par\textit{Protocols such as Raft replicate an ordered log so nodes agree despite failures. Tradeoff: Coordination provides a consistent control plane but costs extra messages and quorum latency. Watch for: Deploying a quorum across unreliable links can reduce availability.}
\paragraph{MCQ-1832 [hard]} What is the central engineering tradeoff of Consensus?
\begin{enumerate}[label=\Alph*.]
\item Headroom handles bursts but idle accelerators are expensive.
\item Coordination provides a consistent control plane but costs extra messages and quorum latency.
\item Success rate improves, while retries amplify overload and tail latency.
\item Stronger guarantees simplify application invariants but increase latency and unavailability.
\end{enumerate}
\noindent\textbf{Answer.} Coordination provides a consistent control plane but costs extra messages and quorum latency.
\par\textit{Protocols such as Raft replicate an ordered log so nodes agree despite failures. Tradeoff: Coordination provides a consistent control plane but costs extra messages and quorum latency. Watch for: Deploying a quorum across unreliable links can reduce availability.}
\paragraph{MCQ-1833 [hard]} Which failure mode should an interviewer expect you to identify for Consensus?
\begin{enumerate}[label=\Alph*.]
\item Deploying a quorum across unreliable links can reduce availability.
\item Planning with average sequence length underestimates tail memory and latency.
\item Layered automatic retries create multiplicative request storms.
\item Quoting CAP without identifying partitions and operations does not guide design.
\end{enumerate}
\noindent\textbf{Answer.} Deploying a quorum across unreliable links can reduce availability.
\par\textit{Protocols such as Raft replicate an ordered log so nodes agree despite failures. Tradeoff: Coordination provides a consistent control plane but costs extra messages and quorum latency. Watch for: Deploying a quorum across unreliable links can reduce availability.}
\paragraph{MCQ-1834 [medium]} A design review is specifically concerned with this risk: Deploying a quorum across unreliable links can reduce availability. Which topic should be reviewed first?
\begin{enumerate}[label=\Alph*.]
\item Consensus
\item CAP and consistency
\item Retries and timeouts
\item Capacity planning
\end{enumerate}
\noindent\textbf{Answer.} Consensus
\par\textit{Protocols such as Raft replicate an ordered log so nodes agree despite failures. Tradeoff: Coordination provides a consistent control plane but costs extra messages and quorum latency. Watch for: Deploying a quorum across unreliable links can reduce availability.}
\paragraph{MCQ-1835 [hard]} Which design note most accurately balances the benefits and costs of Consensus?
\begin{enumerate}[label=\Alph*.]
\item Headroom handles bursts but idle accelerators are expensive.
\item Coordination provides a consistent control plane but costs extra messages and quorum latency.
\item Stronger guarantees simplify application invariants but increase latency and unavailability.
\item Success rate improves, while retries amplify overload and tail latency.
\end{enumerate}
\noindent\textbf{Answer.} Coordination provides a consistent control plane but costs extra messages and quorum latency.
\par\textit{Protocols such as Raft replicate an ordered log so nodes agree despite failures. Tradeoff: Coordination provides a consistent control plane but costs extra messages and quorum latency. Watch for: Deploying a quorum across unreliable links can reduce availability.}
\paragraph{MCQ-1836 [medium]} Which explanation would be strongest in a system-design interview about Consensus?
\begin{enumerate}[label=\Alph*.]
\item Timeouts bound uncertainty; retries use budgets, backoff, jitter, and idempotency for transient failures.
\item Protocols such as Raft replicate an ordered log so nodes agree despite failures.
\item Under a network partition, a distributed system must trade availability against consistent reads and writes.
\item Workload models connect arrival rate, prompt and output tokens, service time, concurrency, memory, and SLO targets.
\end{enumerate}
\noindent\textbf{Answer.} Protocols such as Raft replicate an ordered log so nodes agree despite failures.
\par\textit{Protocols such as Raft replicate an ordered log so nodes agree despite failures. Tradeoff: Coordination provides a consistent control plane but costs extra messages and quorum latency. Watch for: Deploying a quorum across unreliable links can reduce availability.}
\paragraph{MCQ-1837 [hard]} Which production warning is most directly associated with Consensus?
\begin{enumerate}[label=\Alph*.]
\item Planning with average sequence length underestimates tail memory and latency.
\item Layered automatic retries create multiplicative request storms.
\item Deploying a quorum across unreliable links can reduce availability.
\item Quoting CAP without identifying partitions and operations does not guide design.
\end{enumerate}
\noindent\textbf{Answer.} Deploying a quorum across unreliable links can reduce availability.
\par\textit{Protocols such as Raft replicate an ordered log so nodes agree despite failures. Tradeoff: Coordination provides a consistent control plane but costs extra messages and quorum latency. Watch for: Deploying a quorum across unreliable links can reduce availability.}
\paragraph{FC-1838 [medium]} Define Consensus in interview-ready language.
\noindent\textbf{Answer.} Protocols such as Raft replicate an ordered log so nodes agree despite failures.
\par\textit{Protocols such as Raft replicate an ordered log so nodes agree despite failures.}
\paragraph{FC-1839 [hard]} State the main tradeoff and production pitfall for Consensus.
\noindent\textbf{Answer.} Tradeoff: Coordination provides a consistent control plane but costs extra messages and quorum latency. Pitfall: Deploying a quorum across unreliable links can reduce availability.
\par\textit{Tradeoff: Coordination provides a consistent control plane but costs extra messages and quorum latency. Pitfall: Deploying a quorum across unreliable links can reduce availability.}
\paragraph{LA-1840 [hard]} Explain Consensus from first principles, then show how you would evaluate and operate it in production.
\noindent\textbf{Answer.} Start with the mechanism: Protocols such as Raft replicate an ordered log so nodes agree despite failures. Make the tradeoff explicit: Coordination provides a consistent control plane but costs extra messages and quorum latency. Define workload-specific offline metrics and a latency/cost budget, test important slices, instrument the critical path, and create a rollback criterion. The production review must explicitly guard against this failure: Deploying a quorum across unreliable links can reduce availability.
\par\textit{A strong answer connects mechanism, measurable tradeoffs, failure isolation, observability, and rollback rather than listing product features.}
\paragraph{MCQ-1841 [medium]} Which statement best describes Queues and backpressure?
\begin{enumerate}[label=\Alph*.]
\item Balancers route by health, load, locality, affinity, or capacity across replicas.
\item Protocols such as Raft replicate an ordered log so nodes agree despite failures.
\item Bulkheads, cells, quotas, and tenant boundaries prevent one failure or workload from consuming shared resources.
\item Bounded queues absorb bursts while admission control and backpressure prevent producers from overwhelming consumers.
\end{enumerate}
\noindent\textbf{Answer.} Bounded queues absorb bursts while admission control and backpressure prevent producers from overwhelming consumers.
\par\textit{Bounded queues absorb bursts while admission control and backpressure prevent producers from overwhelming consumers. Tradeoff: Smoothing improves utilization, while waiting increases latency and stale work. Watch for: An unbounded queue converts overload into memory exhaustion.}
\paragraph{MCQ-1842 [hard]} What is the central engineering tradeoff of Queues and backpressure?
\begin{enumerate}[label=\Alph*.]
\item Distribution improves availability but generic least-connections may ignore token workload size.
\item Smoothing improves utilization, while waiting increases latency and stale work.
\item Coordination provides a consistent control plane but costs extra messages and quorum latency.
\item Blast radius shrinks at some efficiency cost.
\end{enumerate}
\noindent\textbf{Answer.} Smoothing improves utilization, while waiting increases latency and stale work.
\par\textit{Bounded queues absorb bursts while admission control and backpressure prevent producers from overwhelming consumers. Tradeoff: Smoothing improves utilization, while waiting increases latency and stale work. Watch for: An unbounded queue converts overload into memory exhaustion.}
\paragraph{MCQ-1843 [hard]} Which failure mode should an interviewer expect you to identify for Queues and backpressure?
\begin{enumerate}[label=\Alph*.]
\item A single global vector index or queue can defeat otherwise isolated services.
\item Sending a long-context request to the least-busy replica can still cause head-of-line blocking.
\item Deploying a quorum across unreliable links can reduce availability.
\item An unbounded queue converts overload into memory exhaustion.
\end{enumerate}
\noindent\textbf{Answer.} An unbounded queue converts overload into memory exhaustion.
\par\textit{Bounded queues absorb bursts while admission control and backpressure prevent producers from overwhelming consumers. Tradeoff: Smoothing improves utilization, while waiting increases latency and stale work. Watch for: An unbounded queue converts overload into memory exhaustion.}
\paragraph{MCQ-1844 [medium]} A design review is specifically concerned with this risk: An unbounded queue converts overload into memory exhaustion. Which topic should be reviewed first?
\begin{enumerate}[label=\Alph*.]
\item Load balancing
\item Queues and backpressure
\item Consensus
\item Fault isolation
\end{enumerate}
\noindent\textbf{Answer.} Queues and backpressure
\par\textit{Bounded queues absorb bursts while admission control and backpressure prevent producers from overwhelming consumers. Tradeoff: Smoothing improves utilization, while waiting increases latency and stale work. Watch for: An unbounded queue converts overload into memory exhaustion.}
\paragraph{MCQ-1845 [hard]} Which design note most accurately balances the benefits and costs of Queues and backpressure?
\begin{enumerate}[label=\Alph*.]
\item Distribution improves availability but generic least-connections may ignore token workload size.
\item Smoothing improves utilization, while waiting increases latency and stale work.
\item Blast radius shrinks at some efficiency cost.
\item Coordination provides a consistent control plane but costs extra messages and quorum latency.
\end{enumerate}
\noindent\textbf{Answer.} Smoothing improves utilization, while waiting increases latency and stale work.
\par\textit{Bounded queues absorb bursts while admission control and backpressure prevent producers from overwhelming consumers. Tradeoff: Smoothing improves utilization, while waiting increases latency and stale work. Watch for: An unbounded queue converts overload into memory exhaustion.}
\paragraph{MCQ-1846 [medium]} Which explanation would be strongest in a system-design interview about Queues and backpressure?
\begin{enumerate}[label=\Alph*.]
\item Balancers route by health, load, locality, affinity, or capacity across replicas.
\item Bulkheads, cells, quotas, and tenant boundaries prevent one failure or workload from consuming shared resources.
\item Bounded queues absorb bursts while admission control and backpressure prevent producers from overwhelming consumers.
\item Protocols such as Raft replicate an ordered log so nodes agree despite failures.
\end{enumerate}
\noindent\textbf{Answer.} Bounded queues absorb bursts while admission control and backpressure prevent producers from overwhelming consumers.
\par\textit{Bounded queues absorb bursts while admission control and backpressure prevent producers from overwhelming consumers. Tradeoff: Smoothing improves utilization, while waiting increases latency and stale work. Watch for: An unbounded queue converts overload into memory exhaustion.}
\paragraph{MCQ-1847 [hard]} Which production warning is most directly associated with Queues and backpressure?
\begin{enumerate}[label=\Alph*.]
\item Sending a long-context request to the least-busy replica can still cause head-of-line blocking.
\item Deploying a quorum across unreliable links can reduce availability.
\item A single global vector index or queue can defeat otherwise isolated services.
\item An unbounded queue converts overload into memory exhaustion.
\end{enumerate}
\noindent\textbf{Answer.} An unbounded queue converts overload into memory exhaustion.
\par\textit{Bounded queues absorb bursts while admission control and backpressure prevent producers from overwhelming consumers. Tradeoff: Smoothing improves utilization, while waiting increases latency and stale work. Watch for: An unbounded queue converts overload into memory exhaustion.}
\paragraph{FC-1848 [medium]} Define Queues and backpressure in interview-ready language.
\noindent\textbf{Answer.} Bounded queues absorb bursts while admission control and backpressure prevent producers from overwhelming consumers.
\par\textit{Bounded queues absorb bursts while admission control and backpressure prevent producers from overwhelming consumers.}
\paragraph{FC-1849 [hard]} State the main tradeoff and production pitfall for Queues and backpressure.
\noindent\textbf{Answer.} Tradeoff: Smoothing improves utilization, while waiting increases latency and stale work. Pitfall: An unbounded queue converts overload into memory exhaustion.
\par\textit{Tradeoff: Smoothing improves utilization, while waiting increases latency and stale work. Pitfall: An unbounded queue converts overload into memory exhaustion.}
\paragraph{LA-1850 [hard]} Explain Queues and backpressure from first principles, then show how you would evaluate and operate it in production.
\noindent\textbf{Answer.} Start with the mechanism: Bounded queues absorb bursts while admission control and backpressure prevent producers from overwhelming consumers. Make the tradeoff explicit: Smoothing improves utilization, while waiting increases latency and stale work. Define workload-specific offline metrics and a latency/cost budget, test important slices, instrument the critical path, and create a rollback criterion. The production review must explicitly guard against this failure: An unbounded queue converts overload into memory exhaustion.
\par\textit{A strong answer connects mechanism, measurable tradeoffs, failure isolation, observability, and rollback rather than listing product features.}
\paragraph{MCQ-1851 [medium]} Which statement best describes Retries and timeouts?
\begin{enumerate}[label=\Alph*.]
\item Bulkheads, cells, quotas, and tenant boundaries prevent one failure or workload from consuming shared resources.
\item Timeouts bound uncertainty; retries use budgets, backoff, jitter, and idempotency for transient failures.
\item Under a network partition, a distributed system must trade availability against consistent reads and writes.
\item Workload models connect arrival rate, prompt and output tokens, service time, concurrency, memory, and SLO targets.
\end{enumerate}
\noindent\textbf{Answer.} Timeouts bound uncertainty; retries use budgets, backoff, jitter, and idempotency for transient failures.
\par\textit{Timeouts bound uncertainty; retries use budgets, backoff, jitter, and idempotency for transient failures. Tradeoff: Success rate improves, while retries amplify overload and tail latency. Watch for: Layered automatic retries create multiplicative request storms.}
\paragraph{MCQ-1852 [hard]} What is the central engineering tradeoff of Retries and timeouts?
\begin{enumerate}[label=\Alph*.]
\item Stronger guarantees simplify application invariants but increase latency and unavailability.
\item Success rate improves, while retries amplify overload and tail latency.
\item Blast radius shrinks at some efficiency cost.
\item Headroom handles bursts but idle accelerators are expensive.
\end{enumerate}
\noindent\textbf{Answer.} Success rate improves, while retries amplify overload and tail latency.
\par\textit{Timeouts bound uncertainty; retries use budgets, backoff, jitter, and idempotency for transient failures. Tradeoff: Success rate improves, while retries amplify overload and tail latency. Watch for: Layered automatic retries create multiplicative request storms.}
\paragraph{MCQ-1853 [hard]} Which failure mode should an interviewer expect you to identify for Retries and timeouts?
\begin{enumerate}[label=\Alph*.]
\item A single global vector index or queue can defeat otherwise isolated services.
\item Layered automatic retries create multiplicative request storms.
\item Quoting CAP without identifying partitions and operations does not guide design.
\item Planning with average sequence length underestimates tail memory and latency.
\end{enumerate}
\noindent\textbf{Answer.} Layered automatic retries create multiplicative request storms.
\par\textit{Timeouts bound uncertainty; retries use budgets, backoff, jitter, and idempotency for transient failures. Tradeoff: Success rate improves, while retries amplify overload and tail latency. Watch for: Layered automatic retries create multiplicative request storms.}
\paragraph{MCQ-1854 [medium]} A design review is specifically concerned with this risk: Layered automatic retries create multiplicative request storms. Which topic should be reviewed first?
\begin{enumerate}[label=\Alph*.]
\item CAP and consistency
\item Retries and timeouts
\item Capacity planning
\item Fault isolation
\end{enumerate}
\noindent\textbf{Answer.} Retries and timeouts
\par\textit{Timeouts bound uncertainty; retries use budgets, backoff, jitter, and idempotency for transient failures. Tradeoff: Success rate improves, while retries amplify overload and tail latency. Watch for: Layered automatic retries create multiplicative request storms.}
\paragraph{MCQ-1855 [hard]} Which design note most accurately balances the benefits and costs of Retries and timeouts?
\begin{enumerate}[label=\Alph*.]
\item Stronger guarantees simplify application invariants but increase latency and unavailability.
\item Success rate improves, while retries amplify overload and tail latency.
\item Headroom handles bursts but idle accelerators are expensive.
\item Blast radius shrinks at some efficiency cost.
\end{enumerate}
\noindent\textbf{Answer.} Success rate improves, while retries amplify overload and tail latency.
\par\textit{Timeouts bound uncertainty; retries use budgets, backoff, jitter, and idempotency for transient failures. Tradeoff: Success rate improves, while retries amplify overload and tail latency. Watch for: Layered automatic retries create multiplicative request storms.}
\paragraph{MCQ-1856 [medium]} Which explanation would be strongest in a system-design interview about Retries and timeouts?
\begin{enumerate}[label=\Alph*.]
\item Workload models connect arrival rate, prompt and output tokens, service time, concurrency, memory, and SLO targets.
\item Under a network partition, a distributed system must trade availability against consistent reads and writes.
\item Bulkheads, cells, quotas, and tenant boundaries prevent one failure or workload from consuming shared resources.
\item Timeouts bound uncertainty; retries use budgets, backoff, jitter, and idempotency for transient failures.
\end{enumerate}
\noindent\textbf{Answer.} Timeouts bound uncertainty; retries use budgets, backoff, jitter, and idempotency for transient failures.
\par\textit{Timeouts bound uncertainty; retries use budgets, backoff, jitter, and idempotency for transient failures. Tradeoff: Success rate improves, while retries amplify overload and tail latency. Watch for: Layered automatic retries create multiplicative request storms.}
\paragraph{MCQ-1857 [hard]} Which production warning is most directly associated with Retries and timeouts?
\begin{enumerate}[label=\Alph*.]
\item Planning with average sequence length underestimates tail memory and latency.
\item Quoting CAP without identifying partitions and operations does not guide design.
\item Layered automatic retries create multiplicative request storms.
\item A single global vector index or queue can defeat otherwise isolated services.
\end{enumerate}
\noindent\textbf{Answer.} Layered automatic retries create multiplicative request storms.
\par\textit{Timeouts bound uncertainty; retries use budgets, backoff, jitter, and idempotency for transient failures. Tradeoff: Success rate improves, while retries amplify overload and tail latency. Watch for: Layered automatic retries create multiplicative request storms.}
\paragraph{FC-1858 [medium]} Define Retries and timeouts in interview-ready language.
\noindent\textbf{Answer.} Timeouts bound uncertainty; retries use budgets, backoff, jitter, and idempotency for transient failures.
\par\textit{Timeouts bound uncertainty; retries use budgets, backoff, jitter, and idempotency for transient failures.}
\paragraph{FC-1859 [hard]} State the main tradeoff and production pitfall for Retries and timeouts.
\noindent\textbf{Answer.} Tradeoff: Success rate improves, while retries amplify overload and tail latency. Pitfall: Layered automatic retries create multiplicative request storms.
\par\textit{Tradeoff: Success rate improves, while retries amplify overload and tail latency. Pitfall: Layered automatic retries create multiplicative request storms.}
\paragraph{LA-1860 [hard]} Explain Retries and timeouts from first principles, then show how you would evaluate and operate it in production.
\noindent\textbf{Answer.} Start with the mechanism: Timeouts bound uncertainty; retries use budgets, backoff, jitter, and idempotency for transient failures. Make the tradeoff explicit: Success rate improves, while retries amplify overload and tail latency. Define workload-specific offline metrics and a latency/cost budget, test important slices, instrument the critical path, and create a rollback criterion. The production review must explicitly guard against this failure: Layered automatic retries create multiplicative request storms.
\par\textit{A strong answer connects mechanism, measurable tradeoffs, failure isolation, observability, and rollback rather than listing product features.}
\paragraph{MCQ-1861 [medium]} Which statement best describes Caching?
\begin{enumerate}[label=\Alph*.]
\item Bulkheads, cells, quotas, and tenant boundaries prevent one failure or workload from consuming shared resources.
\item Caches trade staleness and invalidation complexity for lower latency and load.
\item Workload models connect arrival rate, prompt and output tokens, service time, concurrency, memory, and SLO targets.
\item Timeouts bound uncertainty; retries use budgets, backoff, jitter, and idempotency for transient failures.
\end{enumerate}
\noindent\textbf{Answer.} Caches trade staleness and invalidation complexity for lower latency and load.
\par\textit{Caches trade staleness and invalidation complexity for lower latency and load. Tradeoff: Hit rates improve cost, but correctness depends on keys, TTLs, and invalidation. Watch for: Omitting tenant, model, prompt version, or authorization from a cache key leaks data.}
\paragraph{MCQ-1862 [hard]} What is the central engineering tradeoff of Caching?
\begin{enumerate}[label=\Alph*.]
\item Success rate improves, while retries amplify overload and tail latency.
\item Headroom handles bursts but idle accelerators are expensive.
\item Blast radius shrinks at some efficiency cost.
\item Hit rates improve cost, but correctness depends on keys, TTLs, and invalidation.
\end{enumerate}
\noindent\textbf{Answer.} Hit rates improve cost, but correctness depends on keys, TTLs, and invalidation.
\par\textit{Caches trade staleness and invalidation complexity for lower latency and load. Tradeoff: Hit rates improve cost, but correctness depends on keys, TTLs, and invalidation. Watch for: Omitting tenant, model, prompt version, or authorization from a cache key leaks data.}
\paragraph{MCQ-1863 [hard]} Which failure mode should an interviewer expect you to identify for Caching?
\begin{enumerate}[label=\Alph*.]
\item Layered automatic retries create multiplicative request storms.
\item Planning with average sequence length underestimates tail memory and latency.
\item A single global vector index or queue can defeat otherwise isolated services.
\item Omitting tenant, model, prompt version, or authorization from a cache key leaks data.
\end{enumerate}
\noindent\textbf{Answer.} Omitting tenant, model, prompt version, or authorization from a cache key leaks data.
\par\textit{Caches trade staleness and invalidation complexity for lower latency and load. Tradeoff: Hit rates improve cost, but correctness depends on keys, TTLs, and invalidation. Watch for: Omitting tenant, model, prompt version, or authorization from a cache key leaks data.}
\paragraph{MCQ-1864 [medium]} A design review is specifically concerned with this risk: Omitting tenant, model, prompt version, or authorization from a cache key leaks data. Which topic should be reviewed first?
\begin{enumerate}[label=\Alph*.]
\item Capacity planning
\item Caching
\item Retries and timeouts
\item Fault isolation
\end{enumerate}
\noindent\textbf{Answer.} Caching
\par\textit{Caches trade staleness and invalidation complexity for lower latency and load. Tradeoff: Hit rates improve cost, but correctness depends on keys, TTLs, and invalidation. Watch for: Omitting tenant, model, prompt version, or authorization from a cache key leaks data.}
\paragraph{MCQ-1865 [hard]} Which design note most accurately balances the benefits and costs of Caching?
\begin{enumerate}[label=\Alph*.]
\item Hit rates improve cost, but correctness depends on keys, TTLs, and invalidation.
\item Success rate improves, while retries amplify overload and tail latency.
\item Blast radius shrinks at some efficiency cost.
\item Headroom handles bursts but idle accelerators are expensive.
\end{enumerate}
\noindent\textbf{Answer.} Hit rates improve cost, but correctness depends on keys, TTLs, and invalidation.
\par\textit{Caches trade staleness and invalidation complexity for lower latency and load. Tradeoff: Hit rates improve cost, but correctness depends on keys, TTLs, and invalidation. Watch for: Omitting tenant, model, prompt version, or authorization from a cache key leaks data.}
\paragraph{MCQ-1866 [medium]} Which explanation would be strongest in a system-design interview about Caching?
\begin{enumerate}[label=\Alph*.]
\item Timeouts bound uncertainty; retries use budgets, backoff, jitter, and idempotency for transient failures.
\item Caches trade staleness and invalidation complexity for lower latency and load.
\item Workload models connect arrival rate, prompt and output tokens, service time, concurrency, memory, and SLO targets.
\item Bulkheads, cells, quotas, and tenant boundaries prevent one failure or workload from consuming shared resources.
\end{enumerate}
\noindent\textbf{Answer.} Caches trade staleness and invalidation complexity for lower latency and load.
\par\textit{Caches trade staleness and invalidation complexity for lower latency and load. Tradeoff: Hit rates improve cost, but correctness depends on keys, TTLs, and invalidation. Watch for: Omitting tenant, model, prompt version, or authorization from a cache key leaks data.}
\paragraph{MCQ-1867 [hard]} Which production warning is most directly associated with Caching?
\begin{enumerate}[label=\Alph*.]
\item Planning with average sequence length underestimates tail memory and latency.
\item Omitting tenant, model, prompt version, or authorization from a cache key leaks data.
\item Layered automatic retries create multiplicative request storms.
\item A single global vector index or queue can defeat otherwise isolated services.
\end{enumerate}
\noindent\textbf{Answer.} Omitting tenant, model, prompt version, or authorization from a cache key leaks data.
\par\textit{Caches trade staleness and invalidation complexity for lower latency and load. Tradeoff: Hit rates improve cost, but correctness depends on keys, TTLs, and invalidation. Watch for: Omitting tenant, model, prompt version, or authorization from a cache key leaks data.}
\paragraph{FC-1868 [medium]} Define Caching in interview-ready language.
\noindent\textbf{Answer.} Caches trade staleness and invalidation complexity for lower latency and load.
\par\textit{Caches trade staleness and invalidation complexity for lower latency and load.}
\paragraph{FC-1869 [hard]} State the main tradeoff and production pitfall for Caching.
\noindent\textbf{Answer.} Tradeoff: Hit rates improve cost, but correctness depends on keys, TTLs, and invalidation. Pitfall: Omitting tenant, model, prompt version, or authorization from a cache key leaks data.
\par\textit{Tradeoff: Hit rates improve cost, but correctness depends on keys, TTLs, and invalidation. Pitfall: Omitting tenant, model, prompt version, or authorization from a cache key leaks data.}
\paragraph{LA-1870 [hard]} Explain Caching from first principles, then show how you would evaluate and operate it in production.
\noindent\textbf{Answer.} Start with the mechanism: Caches trade staleness and invalidation complexity for lower latency and load. Make the tradeoff explicit: Hit rates improve cost, but correctness depends on keys, TTLs, and invalidation. Define workload-specific offline metrics and a latency/cost budget, test important slices, instrument the critical path, and create a rollback criterion. The production review must explicitly guard against this failure: Omitting tenant, model, prompt version, or authorization from a cache key leaks data.
\par\textit{A strong answer connects mechanism, measurable tradeoffs, failure isolation, observability, and rollback rather than listing product features.}
\paragraph{MCQ-1871 [medium]} Which statement best describes Load balancing?
\begin{enumerate}[label=\Alph*.]
\item Balancers route by health, load, locality, affinity, or capacity across replicas.
\item Bounded queues absorb bursts while admission control and backpressure prevent producers from overwhelming consumers.
\item Workload models connect arrival rate, prompt and output tokens, service time, concurrency, memory, and SLO targets.
\item Systems typically combine at-least-once delivery with idempotent consumers, dedupe keys, and transactional outboxes.
\end{enumerate}
\noindent\textbf{Answer.} Balancers route by health, load, locality, affinity, or capacity across replicas.
\par\textit{Balancers route by health, load, locality, affinity, or capacity across replicas. Tradeoff: Distribution improves availability but generic least-connections may ignore token workload size. Watch for: Sending a long-context request to the least-busy replica can still cause head-of-line blocking.}
\paragraph{MCQ-1872 [hard]} What is the central engineering tradeoff of Load balancing?
\begin{enumerate}[label=\Alph*.]
\item Headroom handles bursts but idle accelerators are expensive.
\item Smoothing improves utilization, while waiting increases latency and stale work.
\item Business effects can be effectively once, but state and retention are needed.
\item Distribution improves availability but generic least-connections may ignore token workload size.
\end{enumerate}
\noindent\textbf{Answer.} Distribution improves availability but generic least-connections may ignore token workload size.
\par\textit{Balancers route by health, load, locality, affinity, or capacity across replicas. Tradeoff: Distribution improves availability but generic least-connections may ignore token workload size. Watch for: Sending a long-context request to the least-busy replica can still cause head-of-line blocking.}
\paragraph{MCQ-1873 [hard]} Which failure mode should an interviewer expect you to identify for Load balancing?
\begin{enumerate}[label=\Alph*.]
\item Planning with average sequence length underestimates tail memory and latency.
\item Equating broker delivery semantics with end-to-end exactly-once effects is incorrect.
\item Sending a long-context request to the least-busy replica can still cause head-of-line blocking.
\item An unbounded queue converts overload into memory exhaustion.
\end{enumerate}
\noindent\textbf{Answer.} Sending a long-context request to the least-busy replica can still cause head-of-line blocking.
\par\textit{Balancers route by health, load, locality, affinity, or capacity across replicas. Tradeoff: Distribution improves availability but generic least-connections may ignore token workload size. Watch for: Sending a long-context request to the least-busy replica can still cause head-of-line blocking.}
\paragraph{MCQ-1874 [medium]} A design review is specifically concerned with this risk: Sending a long-context request to the least-busy replica can still cause head-of-line blocking. Which topic should be reviewed first?
\begin{enumerate}[label=\Alph*.]
\item Queues and backpressure
\item Capacity planning
\item Exactly-once effects
\item Load balancing
\end{enumerate}
\noindent\textbf{Answer.} Load balancing
\par\textit{Balancers route by health, load, locality, affinity, or capacity across replicas. Tradeoff: Distribution improves availability but generic least-connections may ignore token workload size. Watch for: Sending a long-context request to the least-busy replica can still cause head-of-line blocking.}
\paragraph{MCQ-1875 [hard]} Which design note most accurately balances the benefits and costs of Load balancing?
\begin{enumerate}[label=\Alph*.]
\item Headroom handles bursts but idle accelerators are expensive.
\item Business effects can be effectively once, but state and retention are needed.
\item Distribution improves availability but generic least-connections may ignore token workload size.
\item Smoothing improves utilization, while waiting increases latency and stale work.
\end{enumerate}
\noindent\textbf{Answer.} Distribution improves availability but generic least-connections may ignore token workload size.
\par\textit{Balancers route by health, load, locality, affinity, or capacity across replicas. Tradeoff: Distribution improves availability but generic least-connections may ignore token workload size. Watch for: Sending a long-context request to the least-busy replica can still cause head-of-line blocking.}
\paragraph{MCQ-1876 [medium]} Which explanation would be strongest in a system-design interview about Load balancing?
\begin{enumerate}[label=\Alph*.]
\item Systems typically combine at-least-once delivery with idempotent consumers, dedupe keys, and transactional outboxes.
\item Balancers route by health, load, locality, affinity, or capacity across replicas.
\item Workload models connect arrival rate, prompt and output tokens, service time, concurrency, memory, and SLO targets.
\item Bounded queues absorb bursts while admission control and backpressure prevent producers from overwhelming consumers.
\end{enumerate}
\noindent\textbf{Answer.} Balancers route by health, load, locality, affinity, or capacity across replicas.
\par\textit{Balancers route by health, load, locality, affinity, or capacity across replicas. Tradeoff: Distribution improves availability but generic least-connections may ignore token workload size. Watch for: Sending a long-context request to the least-busy replica can still cause head-of-line blocking.}
\paragraph{MCQ-1877 [hard]} Which production warning is most directly associated with Load balancing?
\begin{enumerate}[label=\Alph*.]
\item Sending a long-context request to the least-busy replica can still cause head-of-line blocking.
\item Equating broker delivery semantics with end-to-end exactly-once effects is incorrect.
\item Planning with average sequence length underestimates tail memory and latency.
\item An unbounded queue converts overload into memory exhaustion.
\end{enumerate}
\noindent\textbf{Answer.} Sending a long-context request to the least-busy replica can still cause head-of-line blocking.
\par\textit{Balancers route by health, load, locality, affinity, or capacity across replicas. Tradeoff: Distribution improves availability but generic least-connections may ignore token workload size. Watch for: Sending a long-context request to the least-busy replica can still cause head-of-line blocking.}
\paragraph{FC-1878 [medium]} Define Load balancing in interview-ready language.
\noindent\textbf{Answer.} Balancers route by health, load, locality, affinity, or capacity across replicas.
\par\textit{Balancers route by health, load, locality, affinity, or capacity across replicas.}
\paragraph{FC-1879 [hard]} State the main tradeoff and production pitfall for Load balancing.
\noindent\textbf{Answer.} Tradeoff: Distribution improves availability but generic least-connections may ignore token workload size. Pitfall: Sending a long-context request to the least-busy replica can still cause head-of-line blocking.
\par\textit{Tradeoff: Distribution improves availability but generic least-connections may ignore token workload size. Pitfall: Sending a long-context request to the least-busy replica can still cause head-of-line blocking.}
\paragraph{LA-1880 [hard]} Explain Load balancing from first principles, then show how you would evaluate and operate it in production.
\noindent\textbf{Answer.} Start with the mechanism: Balancers route by health, load, locality, affinity, or capacity across replicas. Make the tradeoff explicit: Distribution improves availability but generic least-connections may ignore token workload size. Define workload-specific offline metrics and a latency/cost budget, test important slices, instrument the critical path, and create a rollback criterion. The production review must explicitly guard against this failure: Sending a long-context request to the least-busy replica can still cause head-of-line blocking.
\par\textit{A strong answer connects mechanism, measurable tradeoffs, failure isolation, observability, and rollback rather than listing product features.}
\paragraph{MCQ-1881 [medium]} Which statement best describes Fault isolation?
\begin{enumerate}[label=\Alph*.]
\item Regions trade user latency, residency, failover speed, replication lag, and operational complexity.
\item Bulkheads, cells, quotas, and tenant boundaries prevent one failure or workload from consuming shared resources.
\item Protocols such as Raft replicate an ordered log so nodes agree despite failures.
\item Balancers route by health, load, locality, affinity, or capacity across replicas.
\end{enumerate}
\noindent\textbf{Answer.} Bulkheads, cells, quotas, and tenant boundaries prevent one failure or workload from consuming shared resources.
\par\textit{Bulkheads, cells, quotas, and tenant boundaries prevent one failure or workload from consuming shared resources. Tradeoff: Blast radius shrinks at some efficiency cost. Watch for: A single global vector index or queue can defeat otherwise isolated services.}
\paragraph{MCQ-1882 [hard]} What is the central engineering tradeoff of Fault isolation?
\begin{enumerate}[label=\Alph*.]
\item Blast radius shrinks at some efficiency cost.
\item Distribution improves availability but generic least-connections may ignore token workload size.
\item Active-active improves availability and latency but makes writes and conflict resolution harder.
\item Coordination provides a consistent control plane but costs extra messages and quorum latency.
\end{enumerate}
\noindent\textbf{Answer.} Blast radius shrinks at some efficiency cost.
\par\textit{Bulkheads, cells, quotas, and tenant boundaries prevent one failure or workload from consuming shared resources. Tradeoff: Blast radius shrinks at some efficiency cost. Watch for: A single global vector index or queue can defeat otherwise isolated services.}
\paragraph{MCQ-1883 [hard]} Which failure mode should an interviewer expect you to identify for Fault isolation?
\begin{enumerate}[label=\Alph*.]
\item Deploying a quorum across unreliable links can reduce availability.
\item A standby region that never receives realistic traffic may fail during disaster recovery.
\item Sending a long-context request to the least-busy replica can still cause head-of-line blocking.
\item A single global vector index or queue can defeat otherwise isolated services.
\end{enumerate}
\noindent\textbf{Answer.} A single global vector index or queue can defeat otherwise isolated services.
\par\textit{Bulkheads, cells, quotas, and tenant boundaries prevent one failure or workload from consuming shared resources. Tradeoff: Blast radius shrinks at some efficiency cost. Watch for: A single global vector index or queue can defeat otherwise isolated services.}
\paragraph{MCQ-1884 [medium]} A design review is specifically concerned with this risk: A single global vector index or queue can defeat otherwise isolated services. Which topic should be reviewed first?
\begin{enumerate}[label=\Alph*.]
\item Load balancing
\item Fault isolation
\item Multi-region design
\item Consensus
\end{enumerate}
\noindent\textbf{Answer.} Fault isolation
\par\textit{Bulkheads, cells, quotas, and tenant boundaries prevent one failure or workload from consuming shared resources. Tradeoff: Blast radius shrinks at some efficiency cost. Watch for: A single global vector index or queue can defeat otherwise isolated services.}
\paragraph{MCQ-1885 [hard]} Which design note most accurately balances the benefits and costs of Fault isolation?
\begin{enumerate}[label=\Alph*.]
\item Coordination provides a consistent control plane but costs extra messages and quorum latency.
\item Blast radius shrinks at some efficiency cost.
\item Active-active improves availability and latency but makes writes and conflict resolution harder.
\item Distribution improves availability but generic least-connections may ignore token workload size.
\end{enumerate}
\noindent\textbf{Answer.} Blast radius shrinks at some efficiency cost.
\par\textit{Bulkheads, cells, quotas, and tenant boundaries prevent one failure or workload from consuming shared resources. Tradeoff: Blast radius shrinks at some efficiency cost. Watch for: A single global vector index or queue can defeat otherwise isolated services.}
\paragraph{MCQ-1886 [medium]} Which explanation would be strongest in a system-design interview about Fault isolation?
\begin{enumerate}[label=\Alph*.]
\item Protocols such as Raft replicate an ordered log so nodes agree despite failures.
\item Regions trade user latency, residency, failover speed, replication lag, and operational complexity.
\item Bulkheads, cells, quotas, and tenant boundaries prevent one failure or workload from consuming shared resources.
\item Balancers route by health, load, locality, affinity, or capacity across replicas.
\end{enumerate}
\noindent\textbf{Answer.} Bulkheads, cells, quotas, and tenant boundaries prevent one failure or workload from consuming shared resources.
\par\textit{Bulkheads, cells, quotas, and tenant boundaries prevent one failure or workload from consuming shared resources. Tradeoff: Blast radius shrinks at some efficiency cost. Watch for: A single global vector index or queue can defeat otherwise isolated services.}
\paragraph{MCQ-1887 [hard]} Which production warning is most directly associated with Fault isolation?
\begin{enumerate}[label=\Alph*.]
\item Deploying a quorum across unreliable links can reduce availability.
\item Sending a long-context request to the least-busy replica can still cause head-of-line blocking.
\item A single global vector index or queue can defeat otherwise isolated services.
\item A standby region that never receives realistic traffic may fail during disaster recovery.
\end{enumerate}
\noindent\textbf{Answer.} A single global vector index or queue can defeat otherwise isolated services.
\par\textit{Bulkheads, cells, quotas, and tenant boundaries prevent one failure or workload from consuming shared resources. Tradeoff: Blast radius shrinks at some efficiency cost. Watch for: A single global vector index or queue can defeat otherwise isolated services.}
\paragraph{FC-1888 [medium]} Define Fault isolation in interview-ready language.
\noindent\textbf{Answer.} Bulkheads, cells, quotas, and tenant boundaries prevent one failure or workload from consuming shared resources.
\par\textit{Bulkheads, cells, quotas, and tenant boundaries prevent one failure or workload from consuming shared resources.}
\paragraph{FC-1889 [hard]} State the main tradeoff and production pitfall for Fault isolation.
\noindent\textbf{Answer.} Tradeoff: Blast radius shrinks at some efficiency cost. Pitfall: A single global vector index or queue can defeat otherwise isolated services.
\par\textit{Tradeoff: Blast radius shrinks at some efficiency cost. Pitfall: A single global vector index or queue can defeat otherwise isolated services.}
\paragraph{LA-1890 [hard]} Explain Fault isolation from first principles, then show how you would evaluate and operate it in production.
\noindent\textbf{Answer.} Start with the mechanism: Bulkheads, cells, quotas, and tenant boundaries prevent one failure or workload from consuming shared resources. Make the tradeoff explicit: Blast radius shrinks at some efficiency cost. Define workload-specific offline metrics and a latency/cost budget, test important slices, instrument the critical path, and create a rollback criterion. The production review must explicitly guard against this failure: A single global vector index or queue can defeat otherwise isolated services.
\par\textit{A strong answer connects mechanism, measurable tradeoffs, failure isolation, observability, and rollback rather than listing product features.}
\paragraph{MCQ-1891 [medium]} Which statement best describes Exactly-once effects?
\begin{enumerate}[label=\Alph*.]
\item Systems typically combine at-least-once delivery with idempotent consumers, dedupe keys, and transactional outboxes.
\item Caches trade staleness and invalidation complexity for lower latency and load.
\item Workload models connect arrival rate, prompt and output tokens, service time, concurrency, memory, and SLO targets.
\item Under a network partition, a distributed system must trade availability against consistent reads and writes.
\end{enumerate}
\noindent\textbf{Answer.} Systems typically combine at-least-once delivery with idempotent consumers, dedupe keys, and transactional outboxes.
\par\textit{Systems typically combine at-least-once delivery with idempotent consumers, dedupe keys, and transactional outboxes. Tradeoff: Business effects can be effectively once, but state and retention are needed. Watch for: Equating broker delivery semantics with end-to-end exactly-once effects is incorrect.}
\paragraph{MCQ-1892 [hard]} What is the central engineering tradeoff of Exactly-once effects?
\begin{enumerate}[label=\Alph*.]
\item Stronger guarantees simplify application invariants but increase latency and unavailability.
\item Hit rates improve cost, but correctness depends on keys, TTLs, and invalidation.
\item Business effects can be effectively once, but state and retention are needed.
\item Headroom handles bursts but idle accelerators are expensive.
\end{enumerate}
\noindent\textbf{Answer.} Business effects can be effectively once, but state and retention are needed.
\par\textit{Systems typically combine at-least-once delivery with idempotent consumers, dedupe keys, and transactional outboxes. Tradeoff: Business effects can be effectively once, but state and retention are needed. Watch for: Equating broker delivery semantics with end-to-end exactly-once effects is incorrect.}
\paragraph{MCQ-1893 [hard]} Which failure mode should an interviewer expect you to identify for Exactly-once effects?
\begin{enumerate}[label=\Alph*.]
\item Quoting CAP without identifying partitions and operations does not guide design.
\item Equating broker delivery semantics with end-to-end exactly-once effects is incorrect.
\item Omitting tenant, model, prompt version, or authorization from a cache key leaks data.
\item Planning with average sequence length underestimates tail memory and latency.
\end{enumerate}
\noindent\textbf{Answer.} Equating broker delivery semantics with end-to-end exactly-once effects is incorrect.
\par\textit{Systems typically combine at-least-once delivery with idempotent consumers, dedupe keys, and transactional outboxes. Tradeoff: Business effects can be effectively once, but state and retention are needed. Watch for: Equating broker delivery semantics with end-to-end exactly-once effects is incorrect.}
\paragraph{MCQ-1894 [medium]} A design review is specifically concerned with this risk: Equating broker delivery semantics with end-to-end exactly-once effects is incorrect. Which topic should be reviewed first?
\begin{enumerate}[label=\Alph*.]
\item Exactly-once effects
\item Capacity planning
\item Caching
\item CAP and consistency
\end{enumerate}
\noindent\textbf{Answer.} Exactly-once effects
\par\textit{Systems typically combine at-least-once delivery with idempotent consumers, dedupe keys, and transactional outboxes. Tradeoff: Business effects can be effectively once, but state and retention are needed. Watch for: Equating broker delivery semantics with end-to-end exactly-once effects is incorrect.}
\paragraph{MCQ-1895 [hard]} Which design note most accurately balances the benefits and costs of Exactly-once effects?
\begin{enumerate}[label=\Alph*.]
\item Business effects can be effectively once, but state and retention are needed.
\item Stronger guarantees simplify application invariants but increase latency and unavailability.
\item Hit rates improve cost, but correctness depends on keys, TTLs, and invalidation.
\item Headroom handles bursts but idle accelerators are expensive.
\end{enumerate}
\noindent\textbf{Answer.} Business effects can be effectively once, but state and retention are needed.
\par\textit{Systems typically combine at-least-once delivery with idempotent consumers, dedupe keys, and transactional outboxes. Tradeoff: Business effects can be effectively once, but state and retention are needed. Watch for: Equating broker delivery semantics with end-to-end exactly-once effects is incorrect.}
\paragraph{MCQ-1896 [medium]} Which explanation would be strongest in a system-design interview about Exactly-once effects?
\begin{enumerate}[label=\Alph*.]
\item Workload models connect arrival rate, prompt and output tokens, service time, concurrency, memory, and SLO targets.
\item Systems typically combine at-least-once delivery with idempotent consumers, dedupe keys, and transactional outboxes.
\item Caches trade staleness and invalidation complexity for lower latency and load.
\item Under a network partition, a distributed system must trade availability against consistent reads and writes.
\end{enumerate}
\noindent\textbf{Answer.} Systems typically combine at-least-once delivery with idempotent consumers, dedupe keys, and transactional outboxes.
\par\textit{Systems typically combine at-least-once delivery with idempotent consumers, dedupe keys, and transactional outboxes. Tradeoff: Business effects can be effectively once, but state and retention are needed. Watch for: Equating broker delivery semantics with end-to-end exactly-once effects is incorrect.}
\paragraph{MCQ-1897 [hard]} Which production warning is most directly associated with Exactly-once effects?
\begin{enumerate}[label=\Alph*.]
\item Equating broker delivery semantics with end-to-end exactly-once effects is incorrect.
\item Omitting tenant, model, prompt version, or authorization from a cache key leaks data.
\item Quoting CAP without identifying partitions and operations does not guide design.
\item Planning with average sequence length underestimates tail memory and latency.
\end{enumerate}
\noindent\textbf{Answer.} Equating broker delivery semantics with end-to-end exactly-once effects is incorrect.
\par\textit{Systems typically combine at-least-once delivery with idempotent consumers, dedupe keys, and transactional outboxes. Tradeoff: Business effects can be effectively once, but state and retention are needed. Watch for: Equating broker delivery semantics with end-to-end exactly-once effects is incorrect.}
\paragraph{FC-1898 [medium]} Define Exactly-once effects in interview-ready language.
\noindent\textbf{Answer.} Systems typically combine at-least-once delivery with idempotent consumers, dedupe keys, and transactional outboxes.
\par\textit{Systems typically combine at-least-once delivery with idempotent consumers, dedupe keys, and transactional outboxes.}
\paragraph{FC-1899 [hard]} State the main tradeoff and production pitfall for Exactly-once effects.
\noindent\textbf{Answer.} Tradeoff: Business effects can be effectively once, but state and retention are needed. Pitfall: Equating broker delivery semantics with end-to-end exactly-once effects is incorrect.
\par\textit{Tradeoff: Business effects can be effectively once, but state and retention are needed. Pitfall: Equating broker delivery semantics with end-to-end exactly-once effects is incorrect.}
\paragraph{LA-1900 [hard]} Explain Exactly-once effects from first principles, then show how you would evaluate and operate it in production.
\noindent\textbf{Answer.} Start with the mechanism: Systems typically combine at-least-once delivery with idempotent consumers, dedupe keys, and transactional outboxes. Make the tradeoff explicit: Business effects can be effectively once, but state and retention are needed. Define workload-specific offline metrics and a latency/cost budget, test important slices, instrument the critical path, and create a rollback criterion. The production review must explicitly guard against this failure: Equating broker delivery semantics with end-to-end exactly-once effects is incorrect.
\par\textit{A strong answer connects mechanism, measurable tradeoffs, failure isolation, observability, and rollback rather than listing product features.}
\paragraph{MCQ-1901 [medium]} Which statement best describes Multi-region design?
\begin{enumerate}[label=\Alph*.]
\item Regions trade user latency, residency, failover speed, replication lag, and operational complexity.
\item Under a network partition, a distributed system must trade availability against consistent reads and writes.
\item Protocols such as Raft replicate an ordered log so nodes agree despite failures.
\item Bounded queues absorb bursts while admission control and backpressure prevent producers from overwhelming consumers.
\end{enumerate}
\noindent\textbf{Answer.} Regions trade user latency, residency, failover speed, replication lag, and operational complexity.
\par\textit{Regions trade user latency, residency, failover speed, replication lag, and operational complexity. Tradeoff: Active-active improves availability and latency but makes writes and conflict resolution harder. Watch for: A standby region that never receives realistic traffic may fail during disaster recovery.}
\paragraph{MCQ-1902 [hard]} What is the central engineering tradeoff of Multi-region design?
\begin{enumerate}[label=\Alph*.]
\item Stronger guarantees simplify application invariants but increase latency and unavailability.
\item Active-active improves availability and latency but makes writes and conflict resolution harder.
\item Coordination provides a consistent control plane but costs extra messages and quorum latency.
\item Smoothing improves utilization, while waiting increases latency and stale work.
\end{enumerate}
\noindent\textbf{Answer.} Active-active improves availability and latency but makes writes and conflict resolution harder.
\par\textit{Regions trade user latency, residency, failover speed, replication lag, and operational complexity. Tradeoff: Active-active improves availability and latency but makes writes and conflict resolution harder. Watch for: A standby region that never receives realistic traffic may fail during disaster recovery.}
\paragraph{MCQ-1903 [hard]} Which failure mode should an interviewer expect you to identify for Multi-region design?
\begin{enumerate}[label=\Alph*.]
\item Deploying a quorum across unreliable links can reduce availability.
\item A standby region that never receives realistic traffic may fail during disaster recovery.
\item Quoting CAP without identifying partitions and operations does not guide design.
\item An unbounded queue converts overload into memory exhaustion.
\end{enumerate}
\noindent\textbf{Answer.} A standby region that never receives realistic traffic may fail during disaster recovery.
\par\textit{Regions trade user latency, residency, failover speed, replication lag, and operational complexity. Tradeoff: Active-active improves availability and latency but makes writes and conflict resolution harder. Watch for: A standby region that never receives realistic traffic may fail during disaster recovery.}
\paragraph{MCQ-1904 [medium]} A design review is specifically concerned with this risk: A standby region that never receives realistic traffic may fail during disaster recovery. Which topic should be reviewed first?
\begin{enumerate}[label=\Alph*.]
\item Consensus
\item Queues and backpressure
\item CAP and consistency
\item Multi-region design
\end{enumerate}
\noindent\textbf{Answer.} Multi-region design
\par\textit{Regions trade user latency, residency, failover speed, replication lag, and operational complexity. Tradeoff: Active-active improves availability and latency but makes writes and conflict resolution harder. Watch for: A standby region that never receives realistic traffic may fail during disaster recovery.}
\paragraph{MCQ-1905 [hard]} Which design note most accurately balances the benefits and costs of Multi-region design?
\begin{enumerate}[label=\Alph*.]
\item Smoothing improves utilization, while waiting increases latency and stale work.
\item Active-active improves availability and latency but makes writes and conflict resolution harder.
\item Coordination provides a consistent control plane but costs extra messages and quorum latency.
\item Stronger guarantees simplify application invariants but increase latency and unavailability.
\end{enumerate}
\noindent\textbf{Answer.} Active-active improves availability and latency but makes writes and conflict resolution harder.
\par\textit{Regions trade user latency, residency, failover speed, replication lag, and operational complexity. Tradeoff: Active-active improves availability and latency but makes writes and conflict resolution harder. Watch for: A standby region that never receives realistic traffic may fail during disaster recovery.}
\paragraph{MCQ-1906 [medium]} Which explanation would be strongest in a system-design interview about Multi-region design?
\begin{enumerate}[label=\Alph*.]
\item Regions trade user latency, residency, failover speed, replication lag, and operational complexity.
\item Under a network partition, a distributed system must trade availability against consistent reads and writes.
\item Protocols such as Raft replicate an ordered log so nodes agree despite failures.
\item Bounded queues absorb bursts while admission control and backpressure prevent producers from overwhelming consumers.
\end{enumerate}
\noindent\textbf{Answer.} Regions trade user latency, residency, failover speed, replication lag, and operational complexity.
\par\textit{Regions trade user latency, residency, failover speed, replication lag, and operational complexity. Tradeoff: Active-active improves availability and latency but makes writes and conflict resolution harder. Watch for: A standby region that never receives realistic traffic may fail during disaster recovery.}
\paragraph{MCQ-1907 [hard]} Which production warning is most directly associated with Multi-region design?
\begin{enumerate}[label=\Alph*.]
\item An unbounded queue converts overload into memory exhaustion.
\item Quoting CAP without identifying partitions and operations does not guide design.
\item A standby region that never receives realistic traffic may fail during disaster recovery.
\item Deploying a quorum across unreliable links can reduce availability.
\end{enumerate}
\noindent\textbf{Answer.} A standby region that never receives realistic traffic may fail during disaster recovery.
\par\textit{Regions trade user latency, residency, failover speed, replication lag, and operational complexity. Tradeoff: Active-active improves availability and latency but makes writes and conflict resolution harder. Watch for: A standby region that never receives realistic traffic may fail during disaster recovery.}
\paragraph{FC-1908 [medium]} Define Multi-region design in interview-ready language.
\noindent\textbf{Answer.} Regions trade user latency, residency, failover speed, replication lag, and operational complexity.
\par\textit{Regions trade user latency, residency, failover speed, replication lag, and operational complexity.}
\paragraph{FC-1909 [hard]} State the main tradeoff and production pitfall for Multi-region design.
\noindent\textbf{Answer.} Tradeoff: Active-active improves availability and latency but makes writes and conflict resolution harder. Pitfall: A standby region that never receives realistic traffic may fail during disaster recovery.
\par\textit{Tradeoff: Active-active improves availability and latency but makes writes and conflict resolution harder. Pitfall: A standby region that never receives realistic traffic may fail during disaster recovery.}
\paragraph{LA-1910 [hard]} Explain Multi-region design from first principles, then show how you would evaluate and operate it in production.
\noindent\textbf{Answer.} Start with the mechanism: Regions trade user latency, residency, failover speed, replication lag, and operational complexity. Make the tradeoff explicit: Active-active improves availability and latency but makes writes and conflict resolution harder. Define workload-specific offline metrics and a latency/cost budget, test important slices, instrument the critical path, and create a rollback criterion. The production review must explicitly guard against this failure: A standby region that never receives realistic traffic may fail during disaster recovery.
\par\textit{A strong answer connects mechanism, measurable tradeoffs, failure isolation, observability, and rollback rather than listing product features.}
\paragraph{MCQ-1911 [medium]} Which statement best describes Capacity planning?
\begin{enumerate}[label=\Alph*.]
\item Systems typically combine at-least-once delivery with idempotent consumers, dedupe keys, and transactional outboxes.
\item Workload models connect arrival rate, prompt and output tokens, service time, concurrency, memory, and SLO targets.
\item Timeouts bound uncertainty; retries use budgets, backoff, jitter, and idempotency for transient failures.
\item Bulkheads, cells, quotas, and tenant boundaries prevent one failure or workload from consuming shared resources.
\end{enumerate}
\noindent\textbf{Answer.} Workload models connect arrival rate, prompt and output tokens, service time, concurrency, memory, and SLO targets.
\par\textit{Workload models connect arrival rate, prompt and output tokens, service time, concurrency, memory, and SLO targets. Tradeoff: Headroom handles bursts but idle accelerators are expensive. Watch for: Planning with average sequence length underestimates tail memory and latency.}
\paragraph{MCQ-1912 [hard]} What is the central engineering tradeoff of Capacity planning?
\begin{enumerate}[label=\Alph*.]
\item Business effects can be effectively once, but state and retention are needed.
\item Success rate improves, while retries amplify overload and tail latency.
\item Headroom handles bursts but idle accelerators are expensive.
\item Blast radius shrinks at some efficiency cost.
\end{enumerate}
\noindent\textbf{Answer.} Headroom handles bursts but idle accelerators are expensive.
\par\textit{Workload models connect arrival rate, prompt and output tokens, service time, concurrency, memory, and SLO targets. Tradeoff: Headroom handles bursts but idle accelerators are expensive. Watch for: Planning with average sequence length underestimates tail memory and latency.}
\paragraph{MCQ-1913 [hard]} Which failure mode should an interviewer expect you to identify for Capacity planning?
\begin{enumerate}[label=\Alph*.]
\item Layered automatic retries create multiplicative request storms.
\item Equating broker delivery semantics with end-to-end exactly-once effects is incorrect.
\item Planning with average sequence length underestimates tail memory and latency.
\item A single global vector index or queue can defeat otherwise isolated services.
\end{enumerate}
\noindent\textbf{Answer.} Planning with average sequence length underestimates tail memory and latency.
\par\textit{Workload models connect arrival rate, prompt and output tokens, service time, concurrency, memory, and SLO targets. Tradeoff: Headroom handles bursts but idle accelerators are expensive. Watch for: Planning with average sequence length underestimates tail memory and latency.}
\paragraph{MCQ-1914 [medium]} A design review is specifically concerned with this risk: Planning with average sequence length underestimates tail memory and latency. Which topic should be reviewed first?
\begin{enumerate}[label=\Alph*.]
\item Capacity planning
\item Retries and timeouts
\item Fault isolation
\item Exactly-once effects
\end{enumerate}
\noindent\textbf{Answer.} Capacity planning
\par\textit{Workload models connect arrival rate, prompt and output tokens, service time, concurrency, memory, and SLO targets. Tradeoff: Headroom handles bursts but idle accelerators are expensive. Watch for: Planning with average sequence length underestimates tail memory and latency.}
\paragraph{MCQ-1915 [hard]} Which design note most accurately balances the benefits and costs of Capacity planning?
\begin{enumerate}[label=\Alph*.]
\item Blast radius shrinks at some efficiency cost.
\item Success rate improves, while retries amplify overload and tail latency.
\item Headroom handles bursts but idle accelerators are expensive.
\item Business effects can be effectively once, but state and retention are needed.
\end{enumerate}
\noindent\textbf{Answer.} Headroom handles bursts but idle accelerators are expensive.
\par\textit{Workload models connect arrival rate, prompt and output tokens, service time, concurrency, memory, and SLO targets. Tradeoff: Headroom handles bursts but idle accelerators are expensive. Watch for: Planning with average sequence length underestimates tail memory and latency.}
\paragraph{MCQ-1916 [medium]} Which explanation would be strongest in a system-design interview about Capacity planning?
\begin{enumerate}[label=\Alph*.]
\item Bulkheads, cells, quotas, and tenant boundaries prevent one failure or workload from consuming shared resources.
\item Timeouts bound uncertainty; retries use budgets, backoff, jitter, and idempotency for transient failures.
\item Systems typically combine at-least-once delivery with idempotent consumers, dedupe keys, and transactional outboxes.
\item Workload models connect arrival rate, prompt and output tokens, service time, concurrency, memory, and SLO targets.
\end{enumerate}
\noindent\textbf{Answer.} Workload models connect arrival rate, prompt and output tokens, service time, concurrency, memory, and SLO targets.
\par\textit{Workload models connect arrival rate, prompt and output tokens, service time, concurrency, memory, and SLO targets. Tradeoff: Headroom handles bursts but idle accelerators are expensive. Watch for: Planning with average sequence length underestimates tail memory and latency.}
\paragraph{MCQ-1917 [hard]} Which production warning is most directly associated with Capacity planning?
\begin{enumerate}[label=\Alph*.]
\item Equating broker delivery semantics with end-to-end exactly-once effects is incorrect.
\item Layered automatic retries create multiplicative request storms.
\item A single global vector index or queue can defeat otherwise isolated services.
\item Planning with average sequence length underestimates tail memory and latency.
\end{enumerate}
\noindent\textbf{Answer.} Planning with average sequence length underestimates tail memory and latency.
\par\textit{Workload models connect arrival rate, prompt and output tokens, service time, concurrency, memory, and SLO targets. Tradeoff: Headroom handles bursts but idle accelerators are expensive. Watch for: Planning with average sequence length underestimates tail memory and latency.}
\paragraph{FC-1918 [medium]} Define Capacity planning in interview-ready language.
\noindent\textbf{Answer.} Workload models connect arrival rate, prompt and output tokens, service time, concurrency, memory, and SLO targets.
\par\textit{Workload models connect arrival rate, prompt and output tokens, service time, concurrency, memory, and SLO targets.}
\paragraph{FC-1919 [hard]} State the main tradeoff and production pitfall for Capacity planning.
\noindent\textbf{Answer.} Tradeoff: Headroom handles bursts but idle accelerators are expensive. Pitfall: Planning with average sequence length underestimates tail memory and latency.
\par\textit{Tradeoff: Headroom handles bursts but idle accelerators are expensive. Pitfall: Planning with average sequence length underestimates tail memory and latency.}
\paragraph{LA-1920 [hard]} Explain Capacity planning from first principles, then show how you would evaluate and operate it in production.
\noindent\textbf{Answer.} Start with the mechanism: Workload models connect arrival rate, prompt and output tokens, service time, concurrency, memory, and SLO targets. Make the tradeoff explicit: Headroom handles bursts but idle accelerators are expensive. Define workload-specific offline metrics and a latency/cost budget, test important slices, instrument the critical path, and create a rollback criterion. The production review must explicitly guard against this failure: Planning with average sequence length underestimates tail memory and latency.
\par\textit{A strong answer connects mechanism, measurable tradeoffs, failure isolation, observability, and rollback rather than listing product features.}
\clearpage\section{Primary sources}
\begin{enumerate}
\item \href{https://arxiv.org/abs/1706.03762}{Attention Is All You Need} (paper)
\item \href{https://arxiv.org/abs/2005.11401}{Retrieval-Augmented Generation for Knowledge-Intensive NLP Tasks} (paper)
\item \href{https://arxiv.org/abs/2004.04906}{Dense Passage Retrieval for Open-Domain Question Answering} (paper)
\item \href{https://arxiv.org/abs/2104.08663}{BEIR: A Heterogeneous Benchmark for Zero-shot Evaluation of Information Retrieval Models} (paper)
\item \href{https://arxiv.org/abs/2309.15217}{RAGAS: Automated Evaluation of Retrieval Augmented Generation} (paper)
\item \href{https://arxiv.org/abs/2109.07958}{TruthfulQA: Measuring How Models Mimic Human Falsehoods} (paper)
\item \href{https://arxiv.org/abs/2210.03629}{ReAct: Synergizing Reasoning and Acting in Language Models} (paper)
\item \href{https://arxiv.org/abs/2303.11366}{Reflexion: Language Agents with Verbal Reinforcement Learning} (paper)
\item \href{https://arxiv.org/abs/2302.04761}{Toolformer: Language Models Can Teach Themselves to Use Tools} (paper)
\item \href{https://arxiv.org/abs/2404.16130}{A Graph RAG Approach to Query-Focused Summarization} (paper)
\item \href{https://arxiv.org/abs/2409.13731}{KAG: Boosting LLMs in Professional Domains via Knowledge Augmented Generation} (paper)
\item \href{https://modelcontextprotocol.io/specification/2026-07-28}{Model Context Protocol Specification (2026-07-28)} (specification)
\item \href{https://arxiv.org/abs/2205.14135}{FlashAttention: Fast and Memory-Efficient Exact Attention with IO-Awareness} (paper)
\item \href{https://arxiv.org/abs/2309.06180}{Efficient Memory Management for Large Language Model Serving with PagedAttention} (paper)
\item \href{https://arxiv.org/abs/2305.13245}{GQA: Training Generalized Multi-Query Transformer Models} (paper)
\item \href{https://arxiv.org/abs/2211.17192}{Fast Inference from Transformers via Speculative Decoding} (paper)
\item \href{https://arxiv.org/abs/2106.09685}{LoRA: Low-Rank Adaptation of Large Language Models} (paper)
\item \href{https://arxiv.org/abs/2305.14314}{QLoRA: Efficient Finetuning of Quantized LLMs} (paper)
\item \href{https://arxiv.org/abs/2305.18290}{Direct Preference Optimization} (paper)
\item \href{https://github.com/pgvector/pgvector}{pgvector Documentation} (official docs)
\item \href{https://docs.pinecone.io/guides/search/hybrid-search}{Pinecone Hybrid Search} (official docs)
\item \href{https://qdrant.tech/documentation/}{Qdrant Documentation} (official docs)
\item \href{https://docs.weaviate.io/weaviate}{Weaviate Documentation} (official docs)
\item \href{https://milvus.io/docs}{Milvus Documentation} (official docs)
\item \href{https://www.elastic.co/docs/solutions/search/vector/knn}{Elasticsearch kNN Search} (official docs)
\item \href{https://neo4j.com/docs/cypher-manual/current/indexes/semantic-indexes/vector-indexes/}{Neo4j Vector Indexes} (official docs)
\item \href{https://docs.langchain.com/oss/python/langgraph/overview}{LangGraph Overview} (official docs)
\item \href{https://docs.temporal.io/}{Temporal Durable Execution} (official docs)
\item \href{https://opentelemetry.io/docs/specs/semconv/}{OpenTelemetry Semantic Conventions} (standard)
\item \href{https://mlflow.org/docs/latest/genai/}{MLflow for LLMs and Agents} (official docs)
\item \href{https://docs.evidentlyai.com/}{Evidently Documentation} (official docs)
\item \href{https://doc.dvc.org/start/data-pipelines/data-pipelines}{DVC Data Pipelines and Versioning} (official docs)
\item \href{https://docs.feast.dev/}{Feast Feature Store Documentation} (official docs)
\item \href{https://docs.vllm.ai/}{vLLM Documentation} (official docs)
\item \href{https://docs.ray.io/en/latest/serve/llm/}{Ray Serve LLM Serving} (official docs)
\item \href{https://kserve.github.io/website/}{KServe Documentation} (official docs)
\item \href{https://kubernetes.io/docs/tasks/run-application/horizontal-pod-autoscale/}{Kubernetes Horizontal Pod Autoscaling} (official docs)
\item \href{https://fastapi.tiangolo.com/}{FastAPI Documentation} (official docs)
\item \href{https://grpc.io/docs/}{gRPC Documentation} (official docs)
\item \href{https://docs.python.org/3/}{Python 3 Documentation} (official docs)
\item \href{https://genai.owasp.org/llm-top-10/}{OWASP Top 10 for LLM Applications} (standard)
\item \href{https://www.nist.gov/itl/ai-risk-management-framework}{NIST AI Risk Management Framework} (standard)
\item \href{https://docs.langchain.com/oss/python/langchain/overview}{LangChain official documentation} (official technology documentation)
\item \href{https://docs.langchain.com/oss/python/langchain/agents}{LangChain agents and runtime} (official technology documentation)
\item \href{https://docs.langchain.com/oss/python/langgraph/graph-api}{LangGraph graph API} (official technology documentation)
\item \href{https://docs.langchain.com/oss/python/langgraph/persistence}{LangGraph persistence design} (official technology documentation)
\item \href{https://docs.langchain.com/oss/python/deepagents/overview}{Deep Agents official documentation} (official technology documentation)
\item \href{https://docs.langchain.com/oss/python/deepagents/harness}{Deep Agents harness} (official technology documentation)
\item \href{https://docs.langchain.com/oss/python/deepagents/backends}{Deep Agents backends} (official technology documentation)
\item \href{https://docs.llamaindex.ai/en/stable/}{LlamaIndex official documentation} (official technology documentation)
\item \href{https://docs.haystack.deepset.ai/docs/intro}{Haystack official documentation} (official technology documentation)
\item \href{https://learn.microsoft.com/en-us/semantic-kernel/overview/}{Semantic Kernel official documentation} (official technology documentation)
\item \href{https://microsoft.github.io/autogen/stable/}{AutoGen official documentation} (official technology documentation)
\item \href{https://docs.crewai.com/}{CrewAI official documentation} (official technology documentation)
\item \href{https://ai.pydantic.dev/}{PydanticAI official documentation} (official technology documentation)
\item \href{https://dspy.ai/}{DSPy official documentation} (official technology documentation)
\item \href{https://openai.github.io/openai-agents-python/}{OpenAI Agents SDK official documentation} (official technology documentation)
\item \href{https://google.github.io/adk-docs/}{Google Agent Development Kit official documentation} (official technology documentation)
\item \href{https://mastra.ai/docs}{Mastra official documentation} (official technology documentation)
\item \href{https://docs.vllm.ai/en/latest/design/arch_overview/}{vLLM architecture} (official technology documentation)
\item \href{https://docs.vllm.ai/en/latest/serving/openai_compatible_server/}{vLLM serving guide} (official technology documentation)
\item \href{https://docs.sglang.ai/}{SGLang official documentation} (official technology documentation)
\item \href{https://huggingface.co/docs/text-generation-inference/index}{Hugging Face Text Generation Inference official documentation} (official technology documentation)
\item \href{https://nvidia.github.io/TensorRT-LLM/}{TensorRT-LLM official documentation} (official technology documentation)
\item \href{https://docs.nvidia.com/deeplearning/triton-inference-server/user-guide/docs/}{NVIDIA Triton Inference Server official documentation} (official technology documentation)
\item \href{https://github.com/ggml-org/llama.cpp}{llama.cpp official documentation} (official technology documentation)
\item \href{https://docs.ollama.com/}{Ollama official documentation} (official technology documentation)
\item \href{https://docs.ray.io/en/latest/serve/}{Ray Serve official documentation} (official technology documentation)
\item \href{https://docs.bentoml.com/}{BentoML official documentation} (official technology documentation)
\item \href{https://localai.io/}{LocalAI official documentation} (official technology documentation)
\item \href{https://modal.com/docs/guide}{Modal official documentation} (official technology documentation)
\item \href{https://docs.baseten.co/}{Baseten official documentation} (official technology documentation)
\item \href{https://huggingface.co/docs/inference-endpoints/index}{Hugging Face Inference Endpoints official documentation} (official technology documentation)
\item \href{https://docs.pinecone.io/}{Pinecone official documentation} (official technology documentation)
\item \href{https://docs.trychroma.com/}{Chroma official documentation} (official technology documentation)
\item \href{https://docs.lancedb.com/}{LanceDB official documentation} (official technology documentation)
\item \href{https://github.com/facebookresearch/faiss}{FAISS official documentation} (official technology documentation)
\item \href{https://www.elastic.co/docs/solutions/search}{Elasticsearch official documentation} (official technology documentation)
\item \href{https://docs.opensearch.org/latest/vector-search/}{OpenSearch official documentation} (official technology documentation)
\item \href{https://docs.vespa.ai/}{Vespa official documentation} (official technology documentation)
\item \href{https://redis.io/docs/latest/develop/ai/}{Redis Search and Vector Sets official documentation} (official technology documentation)
\item \href{https://neo4j.com/docs/}{Neo4j official documentation} (official technology documentation)
\item \href{https://neo4j.com/docs/getting-started/appendix/graphdb-concepts/}{Neo4j graph database concepts} (official technology documentation)
\item \href{https://neo4j.com/docs/cypher-manual/current/}{Cypher manual} (official technology documentation)
\item \href{https://neo4j.com/docs/neo4j-graphrag-python/current/}{Neo4j GraphRAG for Python} (official technology documentation)
\item \href{https://memgraph.com/docs}{Memgraph official documentation} (official technology documentation)
\item \href{https://docs.falkordb.com/}{FalkorDB official documentation} (official technology documentation)
\item \href{https://docs.aws.amazon.com/neptune/}{Amazon Neptune official documentation} (official technology documentation)
\item \href{https://docs.kuzudb.com/}{Kuzu official documentation} (official technology documentation)
\item \href{https://docs.langchain.com/langsmith/home}{LangSmith official documentation} (official technology documentation)
\item \href{https://docs.langchain.com/langsmith/observability}{LangSmith observability} (official technology documentation)
\item \href{https://docs.langchain.com/langsmith/evaluation}{LangSmith evaluation} (official technology documentation)
\item \href{https://langfuse.com/docs}{Langfuse official documentation} (official technology documentation)
\item \href{https://arize.com/docs/phoenix}{Arize Phoenix official documentation} (official technology documentation)
\item \href{https://docs.helicone.ai/}{Helicone official documentation} (official technology documentation)
\item \href{https://opentelemetry.io/docs/}{OpenTelemetry official documentation} (official technology documentation)
\item \href{https://weave-docs.wandb.ai/}{Weights \& Biases Weave official documentation} (official technology documentation)
\item \href{https://www.braintrust.dev/docs}{Braintrust official documentation} (official technology documentation)
\item \href{https://docs.ragas.io/}{Ragas official documentation} (official technology documentation)
\item \href{https://deepeval.com/docs/}{DeepEval official documentation} (official technology documentation)
\item \href{https://www.trulens.org/}{TruLens official documentation} (official technology documentation)
\item \href{https://www.promptfoo.dev/docs/}{Promptfoo official documentation} (official technology documentation)
\item \href{https://modelcontextprotocol.io/docs/learn/architecture}{MCP architecture} (official technology documentation)
\item \href{https://modelcontextprotocol.io/specification/latest/basic/authorization}{MCP authorization} (official technology documentation)
\item \href{https://modelcontextprotocol.io/docs/sdk}{Official MCP SDKs official documentation} (official technology documentation)
\item \href{https://gofastmcp.com/getting-started/welcome}{FastMCP official documentation} (official technology documentation)
\item \href{https://github.com/modelcontextprotocol/inspector}{MCP Inspector official documentation} (official technology documentation)
\item \href{https://docs.docker.com/ai/mcp-catalog-and-toolkit/}{Docker MCP Toolkit and Gateway official documentation} (official technology documentation)
\item \href{https://smithery.ai/docs}{Smithery official documentation} (official technology documentation)
\item \href{https://docs.composio.dev/}{Composio official documentation} (official technology documentation)
\item \href{https://docs.arcade.dev/}{Arcade official documentation} (official technology documentation)
\item \href{https://e2b.dev/docs}{E2B official documentation} (official technology documentation)
\item \href{https://www.daytona.io/docs/}{Daytona official documentation} (official technology documentation)
\item \href{https://modal.com/docs/guide/sandbox}{Modal Sandboxes official documentation} (official technology documentation)
\item \href{https://docs.docker.com/engine/security/}{Docker official documentation} (official technology documentation)
\item \href{https://gvisor.dev/docs/}{gVisor official documentation} (official technology documentation)
\item \href{https://github.com/firecracker-microvm/firecracker}{Firecracker official documentation} (official technology documentation)
\item \href{https://katacontainers.io/docs/}{Kata Containers official documentation} (official technology documentation)
\item \href{https://kubernetes.io/docs/concepts/workloads/controllers/job/}{Kubernetes Jobs official documentation} (official technology documentation)
\item \href{https://airflow.apache.org/docs/}{Apache Airflow official documentation} (official technology documentation)
\item \href{https://docs.prefect.io/}{Prefect official documentation} (official technology documentation)
\item \href{https://docs.dagster.io/}{Dagster official documentation} (official technology documentation)
\item \href{https://dvc.org/doc}{DVC official documentation} (official technology documentation)
\item \href{https://docs.lakefs.io/}{lakeFS official documentation} (official technology documentation)
\item \href{https://www.kubeflow.org/docs/components/pipelines/}{Kubeflow Pipelines official documentation} (official technology documentation)
\item \href{https://docs.zenml.io/}{ZenML official documentation} (official technology documentation)
\item \href{https://help.getzep.com/}{Zep official documentation} (official technology documentation)
\item \href{https://docs.mem0.ai/}{Mem0 official documentation} (official technology documentation)
\item \href{https://docs.pytorch.org/docs/stable/index.html}{PyTorch official documentation} (official technology documentation)
\item \href{https://docs.jax.dev/en/latest/}{JAX official documentation} (official technology documentation)
\item \href{https://huggingface.co/docs/transformers/index}{Hugging Face Transformers official documentation} (official technology documentation)
\item \href{https://sbert.net/}{Sentence Transformers official documentation} (official technology documentation)
\item \href{https://huggingface.co/docs/peft/index}{Hugging Face PEFT official documentation} (official technology documentation)
\item \href{https://huggingface.co/docs/bitsandbytes/main/en/index}{bitsandbytes official documentation} (official technology documentation)
\item \href{https://onnxruntime.ai/docs/}{ONNX Runtime official documentation} (official technology documentation)
\item \href{https://docs.openvino.ai/}{OpenVINO official documentation} (official technology documentation)
\item \href{https://bge-model.com/}{FlagEmbedding official documentation} (official technology documentation)
\item \href{https://python.useinstructor.com/}{Instructor official documentation} (official technology documentation)
\item \href{https://docs.cohere.com/docs/rerank-overview}{Cohere Rerank official documentation} (official technology documentation)
\item \href{https://huggingface.co/docs/text-embeddings-inference/index}{Hugging Face Text Embeddings Inference official documentation} (official technology documentation)
\item \href{https://docs.litellm.ai/}{LiteLLM official documentation} (official technology documentation)
\item \href{https://portkey.ai/docs/}{Portkey AI Gateway official documentation} (official technology documentation)
\item \href{https://developer.konghq.com/ai-gateway/}{Kong AI Gateway official documentation} (official technology documentation)
\item \href{https://aigateway.envoyproxy.io/docs/}{Envoy AI Gateway official documentation} (official technology documentation)
\item \href{https://docs.aws.amazon.com/sagemaker/latest/dg/deploy-model.html}{Amazon SageMaker AI Endpoints official documentation} (official technology documentation)
\item \href{https://cloud.google.com/vertex-ai/docs/predictions/overview}{Vertex AI Endpoints official documentation} (official technology documentation)
\item \href{https://learn.microsoft.com/en-us/azure/machine-learning/concept-endpoints-online}{Azure Machine Learning Managed Online Endpoints official documentation} (official technology documentation)
\item \href{https://docs.nvidia.com/nemo/guardrails/latest/}{NVIDIA NeMo Guardrails official documentation} (official technology documentation)
\item \href{https://www.guardrailsai.com/docs}{Guardrails AI official documentation} (official technology documentation)
\item \href{https://www.llama.com/docs/model-cards-and-prompt-formats/llama-guard-4/}{Llama Guard official documentation} (official technology documentation)
\item \href{https://protectai.github.io/llm-guard/}{Protect AI LLM Guard official documentation} (official technology documentation)
\item \href{https://microsoft.github.io/presidio/}{Microsoft Presidio official documentation} (official technology documentation)
\item \href{https://www.openpolicyagent.org/docs/latest/}{Open Policy Agent official documentation} (official technology documentation)
\item \href{https://docs.cedarpolicy.com/}{Cedar official documentation} (official technology documentation)
\item \href{https://docs.lakera.ai/docs}{Lakera Guard official documentation} (official technology documentation)
\item \href{https://docs.agno.com/}{Agno official documentation} (official technology documentation)
\item \href{https://huggingface.co/docs/smolagents/index}{Hugging Face smolagents official documentation} (official technology documentation)
\item \href{https://docs.letta.com/}{Letta official documentation} (official technology documentation)
\item \href{https://langroid.github.io/langroid/}{Langroid official documentation} (official technology documentation)
\item \href{https://framework.beeai.dev/}{BeeAI Framework official documentation} (official technology documentation)
\item \href{https://github.com/strands-agents/harness-sdk}{Strands Agents SDK official documentation} (official technology documentation)
\item \href{https://platform.claude.com/docs/en/agent-sdk/overview}{Claude Agent SDK official documentation} (official technology documentation)
\item \href{https://ai-sdk.dev/docs/introduction}{Vercel AI SDK official documentation} (official technology documentation)
\item \href{https://learn.microsoft.com/en-us/agent-framework/overview/}{Microsoft Agent Framework official documentation} (official technology documentation)
\item \href{https://docs.camel-ai.org/}{CAMEL-AI official documentation} (official technology documentation)
\item \href{https://rasa.com/docs/}{Rasa official documentation} (official technology documentation)
\item \href{https://docs.langflow.org/}{Langflow official documentation} (official technology documentation)
\item \href{https://docs.dify.ai/}{Dify official documentation} (official technology documentation)
\item \href{https://docs.flowiseai.com/}{Flowise official documentation} (official technology documentation)
\item \href{https://developers.cloudflare.com/agents/}{Cloudflare Agents SDK official documentation} (official technology documentation)
\item \href{https://huggingface.co/docs/accelerate/index}{Hugging Face Accelerate official documentation} (official technology documentation)
\item \href{https://www.deepspeed.ai/}{DeepSpeed official documentation} (official technology documentation)
\item \href{https://huggingface.co/docs/trl/index}{Hugging Face TRL official documentation} (official technology documentation)
\item \href{https://docs.axolotl.ai/}{Axolotl official documentation} (official technology documentation)
\item \href{https://docs.unsloth.ai/}{Unsloth official documentation} (official technology documentation)
\item \href{https://llm-d.ai/docs/dev}{llm-d official documentation} (official technology documentation)
\item \href{https://llm-d.ai/docs/architecture}{llm-d architecture} (official technology documentation)
\item \href{https://kserve.github.io/website/docs/model-serving/generative-inference/llmisvc/llmisvc-overview}{KServe LLMInferenceService} (official technology documentation)
\item \href{https://docs.nvidia.com/dynamo/latest/}{NVIDIA Dynamo official documentation} (official technology documentation)
\item \href{https://docs.nvidia.com/nim/large-language-models/latest/introduction.html}{NVIDIA NIM official documentation} (official technology documentation)
\item \href{https://docs.nvidia.com/nim/large-language-models/latest/reference/architecture.html}{NIM LLM architecture} (official technology documentation)
\item \href{https://docs.pytorch.org/serve/}{TorchServe official documentation} (official technology documentation)
\item \href{https://deepspeed-mii.readthedocs.io/en/latest/}{DeepSpeed-MII official documentation} (official technology documentation)
\item \href{https://github.com/bentoml/OpenLLM}{OpenLLM official documentation} (official technology documentation)
\item \href{https://lmdeploy.readthedocs.io/en/latest/}{LMDeploy official documentation} (official technology documentation)
\item \href{https://llm.mlc.ai/docs/}{MLC LLM official documentation} (official technology documentation)
\item \href{https://replicate.com/docs}{Replicate official documentation} (official technology documentation)
\item \href{https://docs.runpod.io/serverless/overview}{Runpod Serverless official documentation} (official technology documentation)
\item \href{https://docs.together.ai/docs/inference-overview}{Together AI Inference official documentation} (official technology documentation)
\item \href{https://docs.fireworks.ai/}{Fireworks AI official documentation} (official technology documentation)
\item \href{https://console.groq.com/docs/overview}{GroqCloud official documentation} (official technology documentation)
\item \href{https://inference-docs.cerebras.ai/introduction}{Cerebras Inference official documentation} (official technology documentation)
\item \href{https://docs.aws.amazon.com/bedrock/latest/userguide/what-is-bedrock.html}{Amazon Bedrock official documentation} (official technology documentation)
\item \href{https://developers.cloudflare.com/workers-ai/}{Cloudflare Workers AI official documentation} (official technology documentation)
\item \href{https://docs.databricks.com/aws/en/machine-learning/model-serving/}{Databricks Model Serving official documentation} (official technology documentation)
\item \href{https://www.mongodb.com/docs/atlas/atlas-vector-search/}{MongoDB Atlas Vector Search official documentation} (official technology documentation)
\item \href{https://docs.couchbase.com/server/current/vector-search/vector-search.html}{Couchbase Vector Search official documentation} (official technology documentation)
\item \href{https://docs.databricks.com/aws/en/generative-ai/vector-search/}{Databricks Mosaic AI Vector Search official documentation} (official technology documentation)
\item \href{https://docs.snowflake.com/en/user-guide/snowflake-cortex/cortex-search/cortex-search-overview}{Snowflake Cortex Search official documentation} (official technology documentation)
\item \href{https://docs.oracle.com/en/database/oracle/oracle-database/23/vecse/}{Oracle AI Vector Search official documentation} (official technology documentation)
\item \href{https://docs.tigergraph.com/}{TigerGraph official documentation} (official technology documentation)
\item \href{https://docs.arangodb.com/}{ArangoDB official documentation} (official technology documentation)
\item \href{https://age.apache.org/age-manual/master/intro/overview.html}{Apache AGE official documentation} (official technology documentation)
\item \href{https://docs.janusgraph.org/}{JanusGraph official documentation} (official technology documentation)
\item \href{https://www.comet.com/docs/opik/}{Opik official documentation} (official technology documentation)
\item \href{https://docs.agentops.ai/}{AgentOps official documentation} (official technology documentation)
\item \href{https://docs.parea.ai/}{Parea official documentation} (official technology documentation)
\item \href{https://docs.honeyhive.ai/}{HoneyHive official documentation} (official technology documentation)
\item \href{https://docs.galileo.ai/}{Galileo official documentation} (official technology documentation)
\item \href{https://docs.giskard.ai/}{Giskard official documentation} (official technology documentation)
\item \href{https://inspect.aisi.org.uk/}{Inspect AI official documentation} (official technology documentation)
\item \href{https://www.getmaxim.ai/docs/}{Maxim AI official documentation} (official technology documentation)
\item \href{https://reference.garak.ai/en/latest/}{garak official documentation} (official technology documentation)
\item \href{https://azure.github.io/PyRIT/}{PyRIT official documentation} (official technology documentation)
\item \href{https://github.com/protectai/rebuff}{Rebuff official documentation} (official technology documentation)
\item \href{https://docs.aws.amazon.com/bedrock/latest/userguide/guardrails.html}{Amazon Bedrock Guardrails official documentation} (official technology documentation)
\item \href{https://learn.microsoft.com/en-us/azure/ai-services/content-safety/overview}{Azure AI Content Safety official documentation} (official technology documentation)
\item \href{https://cloud.google.com/security-command-center/docs/model-armor-overview}{Google Cloud Model Armor official documentation} (official technology documentation)
\item \href{https://modelcontextprotocol.io/registry/about}{Official MCP Registry official documentation} (official technology documentation)
\item \href{https://registry.modelcontextprotocol.io/docs}{Registry API} (official technology documentation)
\item \href{https://github.com/github/github-mcp-server}{GitHub MCP Server official documentation} (official technology documentation)
\item \href{https://github.com/awslabs/mcp}{AWS MCP Servers official documentation} (official technology documentation)
\item \href{https://googleapis.github.io/genai-toolbox/}{MCP Toolbox for Databases official documentation} (official technology documentation)
\item \href{https://github.com/microsoft/mcp}{Azure MCP Server official documentation} (official technology documentation)
\item \href{https://docs.zapier.com/mcp/home}{Zapier MCP official documentation} (official technology documentation)
\item \href{https://mcp-use.com/docs/}{mcp-use official documentation} (official technology documentation)
\item \href{https://developers.cloudflare.com/sandbox/}{Cloudflare Sandbox SDK official documentation} (official technology documentation)
\item \href{https://developers.cloudflare.com/sandbox/api/storage/}{Sandbox storage and persistence} (official technology documentation)
\item \href{https://docs.runloop.ai/}{Runloop official documentation} (official technology documentation)
\item \href{https://docs.sprites.dev/}{Sprites official documentation} (official technology documentation)
\item \href{https://docs.browserbase.com/}{Browserbase official documentation} (official technology documentation)
\item \href{https://docs.docker.com/ai/sandboxes/}{Docker Sandboxes official documentation} (official technology documentation)
\item \href{https://docs.wasmtime.dev/}{Wasmtime official documentation} (official technology documentation)
\item \href{https://github.com/google/nsjail}{nsjail official documentation} (official technology documentation)
\item \href{https://argo-workflows.readthedocs.io/}{Argo Workflows official documentation} (official technology documentation)
\item \href{https://docs.flyte.org/}{Flyte official documentation} (official technology documentation)
\item \href{https://docs.metaflow.org/}{Metaflow official documentation} (official technology documentation)
\item \href{https://clear.ml/docs/latest/docs/}{ClearML official documentation} (official technology documentation)
\item \href{https://labelstud.io/guide/}{Label Studio official documentation} (official technology documentation)
\item \href{https://docs.greatexpectations.io/docs/}{Great Expectations official documentation} (official technology documentation)
\item \href{https://docs.airbyte.com/}{Airbyte official documentation} (official technology documentation)
\item \href{https://docs.unstructured.io/}{Unstructured official documentation} (official technology documentation)
\item \href{https://docling-project.github.io/docling/}{Docling official documentation} (official technology documentation)
\end{enumerate}
